% Studies-in-Algebra-and-Group-Theory - Lecture 7: More maps and spaces
% Author: S. D. V. Stephens
% Date: 2025-12-21
% Compile with: pdflatex -shell-escape

\documentclass{report}
\input{~/university/preamble.tex}
% preamble-darkmode.tex
% Dark mode extension for lecture notes
% Usage: % preamble-darkmode.tex
% Dark mode extension for lecture notes
% Usage: % preamble-darkmode.tex
% Dark mode extension for lecture notes
% Usage: \input{~/university/preamble-darkmode.tex} AFTER preamble.tex
%
% This provides:
% - Dark background with light text
% - Material Design color palette
% - Ocean blue gradient colors (p1-p9)
% - Enhanced TikZ/PGFPlots for dark backgrounds
% - qq tcolorbox environment
% - tkz-euclide for geometric constructions

% ===========================================
% DARK MODE PACKAGE
% ===========================================
\usepackage{darkmode}
\enabledarkmode

% ===========================================
% ADDITIONAL PACKAGES
% ===========================================
\usepackage{changepage}
\usepackage{ulem}
\usepackage{colortbl}
\usepackage{tkz-euclide}
\usepackage{capt-of}

% ===========================================
% EXTENDED TIKZ LIBRARIES
% ===========================================
\usetikzlibrary{patterns,3d,backgrounds}
\usepgfplotslibrary{groupplots}

% Upgrade pgfplots compatibility
\pgfplotsset{compat=1.18}

% ===========================================
% OCEAN BLUE GRADIENT (p1-p9)
% ===========================================
\definecolor{p1}{HTML}{caf0f8}
\definecolor{p2}{HTML}{ade8f4}
\definecolor{p3}{HTML}{90e0ef}
\definecolor{p4}{HTML}{48cae4}
\definecolor{p5}{HTML}{00b4d8}
\definecolor{p6}{HTML}{0096c7}
\definecolor{p7}{HTML}{0077b6}
\definecolor{p8}{HTML}{023e8a}
\definecolor{p9}{HTML}{03045e}

% Dark background color
\definecolor{pag}{HTML}{293133}

% ===========================================
% MATERIAL DESIGN COLORS
% ===========================================
% Note: myred, myyellow already defined in main preamble
% These are additional/alternate versions
\definecolor{matred}{HTML}{f44336}
\definecolor{matpink}{HTML}{e81e63}
\definecolor{matpurple}{HTML}{9c27b0}
\definecolor{matdeeppurple}{HTML}{673ab7}
\definecolor{matindigo}{HTML}{3f51b5}
\definecolor{matblue}{HTML}{2196f3}
\definecolor{matlightblue}{HTML}{03a9f4}
\definecolor{matcyan}{HTML}{00bcd4}
\definecolor{matteal}{HTML}{009688}
\definecolor{matgreen}{HTML}{4caf50}
\definecolor{matlightgreen}{HTML}{8bc34a}
\definecolor{matlime}{HTML}{cddc39}
\definecolor{matyellow}{HTML}{ffeb3b}
\definecolor{matamber}{HTML}{ffc107}
\definecolor{matorange}{HTML}{ff9800}
\definecolor{matdeeporange}{HTML}{ff5722}
\definecolor{matbrown}{HTML}{795548}
\definecolor{matgray}{HTML}{9e9e9e}
\definecolor{matbluegray}{HTML}{607d8b}

% ===========================================
% COLORLET FOR TIKZ
% ===========================================
\colorlet{xcol}{blue!60!black}

% ===========================================
% DARK MODE TCOLORBOX
% ===========================================
\newtcolorbox{qq}[2][]{%
    boxrule=0.75pt,
    sharp corners,
    colframe=white,
    colback=pag,
    coltext=white,
    #1,
}

% Dark definition box
\newtcolorbox{darkdfn}[2][]{%
    enhanced,
    boxrule=1pt,
    sharp corners,
    colframe=p5,
    colback=pag,
    coltext=white,
    coltitle=white,
    fonttitle=\bfseries,
    title=#2,
    #1,
}

% Dark theorem box
\newtcolorbox{darkthm}[2][]{%
    enhanced,
    boxrule=1pt,
    sharp corners,
    colframe=matpurple,
    colback=pag,
    coltext=white,
    coltitle=white,
    fonttitle=\bfseries,
    title=#2,
    #1,
}

% ===========================================
% TIKZ 3D FIX
% ===========================================
% Small fix for canvas is xy plane at z
% https://tex.stackexchange.com/a/48776/121799
\makeatletter
\tikzoption{canvas is xy plane at z}[]{%
    \def\tikz@plane@origin{\pgfpointxyz{0}{0}{#1}}%
    \def\tikz@plane@x{\pgfpointxyz{1}{0}{#1}}%
    \def\tikz@plane@y{\pgfpointxyz{0}{1}{#1}}%
    \tikz@canvas@is@plane}
\makeatother

% ===========================================
% CONDITIONAL LINE TIKZSET
% ===========================================
\tikzset{
    conditional line/.style 2 args={
        /utils/exec={\pgfmathsetmacro{\xcoordA}{#1}},
        /utils/exec={\pgfmathsetmacro{\xcoordB}{#2}},
        insert path={
            \ifdim\xcoordA pt>\xcoordB pt
            (\xcoordA,0) -- (\xcoordB,0)
            \fi
        },
    },
}

% ===========================================
% PGFPLOTS DARK MODE DEFAULTS
% ===========================================
\pgfplotsset{
    dark axis/.style={
        axis line style={white},
        tick style={white},
        label style={white},
        grid style={pag!70},
        every axis label/.append style={white},
        every tick label/.append style={white},
    },
}

% ===========================================
% TIKZ EXTERNALIZATION (for fast compilation)
% ===========================================
% This creates .md5 cache files for figures
% Requires: pdflatex -shell-escape
%
% To enable, add AFTER loading this file:
%   \enabletikzexternalize
%
\newcommand{\enabletikzexternalize}{%
  \usetikzlibrary{external}%
  \tikzexternalize[prefix=figures/]%
}

% ===========================================
% CONVENIENT SHORTCUTS FOR DARK PLOTS
% ===========================================
\newcommand{\darkplot}[1]{%
    \begin{tikzpicture}
    \begin{axis}[
        dark axis,
        grid=both,
        major grid style={line width=.2pt,draw=pag!70},
        #1
    ]
}


% ===========================================
% QUICK GEOMETRY MACROS (tkz-euclide)
% ===========================================
% Draw a point: \point{name}{x}{y}
\newcommand{\gpoint}[3]{\tkzDefPoint(#2,#3){#1}}

% Draw line through points: \gline{A}{B}
\newcommand{\gline}[2]{\tkzDrawLine[color=p4](#1,#2)}

% Draw circle: \gcircle{center}{radius}
\newcommand{\gcircle}[2]{\tkzDrawCircle[color=p5](#1,#2)}

% Mark angle: \gangle{A}{B}{C}
\newcommand{\gangle}[3]{\tkzMarkAngle[color=p3,size=0.5](#1,#2,#3)}

% ===========================================
% USAGE EXAMPLE
% ===========================================
% In your document:
%
% \begin{tikzpicture}
%   \begin{axis}[dark axis, ...]
%     \addplot[p4, thick] {sin(deg(x))};
%   \end{axis}
% \end{tikzpicture}
%
% Or with qq box:
% \begin{qq}{My Title}
%   Content here in dark box
% \end{qq}
 AFTER preamble.tex
%
% This provides:
% - Dark background with light text
% - Material Design color palette
% - Ocean blue gradient colors (p1-p9)
% - Enhanced TikZ/PGFPlots for dark backgrounds
% - qq tcolorbox environment
% - tkz-euclide for geometric constructions

% ===========================================
% DARK MODE PACKAGE
% ===========================================
\usepackage{darkmode}
\enabledarkmode

% ===========================================
% ADDITIONAL PACKAGES
% ===========================================
\usepackage{changepage}
\usepackage{ulem}
\usepackage{colortbl}
\usepackage{tkz-euclide}
\usepackage{capt-of}

% ===========================================
% EXTENDED TIKZ LIBRARIES
% ===========================================
\usetikzlibrary{patterns,3d,backgrounds}
\usepgfplotslibrary{groupplots}

% Upgrade pgfplots compatibility
\pgfplotsset{compat=1.18}

% ===========================================
% OCEAN BLUE GRADIENT (p1-p9)
% ===========================================
\definecolor{p1}{HTML}{caf0f8}
\definecolor{p2}{HTML}{ade8f4}
\definecolor{p3}{HTML}{90e0ef}
\definecolor{p4}{HTML}{48cae4}
\definecolor{p5}{HTML}{00b4d8}
\definecolor{p6}{HTML}{0096c7}
\definecolor{p7}{HTML}{0077b6}
\definecolor{p8}{HTML}{023e8a}
\definecolor{p9}{HTML}{03045e}

% Dark background color
\definecolor{pag}{HTML}{293133}

% ===========================================
% MATERIAL DESIGN COLORS
% ===========================================
% Note: myred, myyellow already defined in main preamble
% These are additional/alternate versions
\definecolor{matred}{HTML}{f44336}
\definecolor{matpink}{HTML}{e81e63}
\definecolor{matpurple}{HTML}{9c27b0}
\definecolor{matdeeppurple}{HTML}{673ab7}
\definecolor{matindigo}{HTML}{3f51b5}
\definecolor{matblue}{HTML}{2196f3}
\definecolor{matlightblue}{HTML}{03a9f4}
\definecolor{matcyan}{HTML}{00bcd4}
\definecolor{matteal}{HTML}{009688}
\definecolor{matgreen}{HTML}{4caf50}
\definecolor{matlightgreen}{HTML}{8bc34a}
\definecolor{matlime}{HTML}{cddc39}
\definecolor{matyellow}{HTML}{ffeb3b}
\definecolor{matamber}{HTML}{ffc107}
\definecolor{matorange}{HTML}{ff9800}
\definecolor{matdeeporange}{HTML}{ff5722}
\definecolor{matbrown}{HTML}{795548}
\definecolor{matgray}{HTML}{9e9e9e}
\definecolor{matbluegray}{HTML}{607d8b}

% ===========================================
% COLORLET FOR TIKZ
% ===========================================
\colorlet{xcol}{blue!60!black}

% ===========================================
% DARK MODE TCOLORBOX
% ===========================================
\newtcolorbox{qq}[2][]{%
    boxrule=0.75pt,
    sharp corners,
    colframe=white,
    colback=pag,
    coltext=white,
    #1,
}

% Dark definition box
\newtcolorbox{darkdfn}[2][]{%
    enhanced,
    boxrule=1pt,
    sharp corners,
    colframe=p5,
    colback=pag,
    coltext=white,
    coltitle=white,
    fonttitle=\bfseries,
    title=#2,
    #1,
}

% Dark theorem box
\newtcolorbox{darkthm}[2][]{%
    enhanced,
    boxrule=1pt,
    sharp corners,
    colframe=matpurple,
    colback=pag,
    coltext=white,
    coltitle=white,
    fonttitle=\bfseries,
    title=#2,
    #1,
}

% ===========================================
% TIKZ 3D FIX
% ===========================================
% Small fix for canvas is xy plane at z
% https://tex.stackexchange.com/a/48776/121799
\makeatletter
\tikzoption{canvas is xy plane at z}[]{%
    \def\tikz@plane@origin{\pgfpointxyz{0}{0}{#1}}%
    \def\tikz@plane@x{\pgfpointxyz{1}{0}{#1}}%
    \def\tikz@plane@y{\pgfpointxyz{0}{1}{#1}}%
    \tikz@canvas@is@plane}
\makeatother

% ===========================================
% CONDITIONAL LINE TIKZSET
% ===========================================
\tikzset{
    conditional line/.style 2 args={
        /utils/exec={\pgfmathsetmacro{\xcoordA}{#1}},
        /utils/exec={\pgfmathsetmacro{\xcoordB}{#2}},
        insert path={
            \ifdim\xcoordA pt>\xcoordB pt
            (\xcoordA,0) -- (\xcoordB,0)
            \fi
        },
    },
}

% ===========================================
% PGFPLOTS DARK MODE DEFAULTS
% ===========================================
\pgfplotsset{
    dark axis/.style={
        axis line style={white},
        tick style={white},
        label style={white},
        grid style={pag!70},
        every axis label/.append style={white},
        every tick label/.append style={white},
    },
}

% ===========================================
% TIKZ EXTERNALIZATION (for fast compilation)
% ===========================================
% This creates .md5 cache files for figures
% Requires: pdflatex -shell-escape
%
% To enable, add AFTER loading this file:
%   \enabletikzexternalize
%
\newcommand{\enabletikzexternalize}{%
  \usetikzlibrary{external}%
  \tikzexternalize[prefix=figures/]%
}

% ===========================================
% CONVENIENT SHORTCUTS FOR DARK PLOTS
% ===========================================
\newcommand{\darkplot}[1]{%
    \begin{tikzpicture}
    \begin{axis}[
        dark axis,
        grid=both,
        major grid style={line width=.2pt,draw=pag!70},
        #1
    ]
}


% ===========================================
% QUICK GEOMETRY MACROS (tkz-euclide)
% ===========================================
% Draw a point: \point{name}{x}{y}
\newcommand{\gpoint}[3]{\tkzDefPoint(#2,#3){#1}}

% Draw line through points: \gline{A}{B}
\newcommand{\gline}[2]{\tkzDrawLine[color=p4](#1,#2)}

% Draw circle: \gcircle{center}{radius}
\newcommand{\gcircle}[2]{\tkzDrawCircle[color=p5](#1,#2)}

% Mark angle: \gangle{A}{B}{C}
\newcommand{\gangle}[3]{\tkzMarkAngle[color=p3,size=0.5](#1,#2,#3)}

% ===========================================
% USAGE EXAMPLE
% ===========================================
% In your document:
%
% \begin{tikzpicture}
%   \begin{axis}[dark axis, ...]
%     \addplot[p4, thick] {sin(deg(x))};
%   \end{axis}
% \end{tikzpicture}
%
% Or with qq box:
% \begin{qq}{My Title}
%   Content here in dark box
% \end{qq}
 AFTER preamble.tex
%
% This provides:
% - Dark background with light text
% - Material Design color palette
% - Ocean blue gradient colors (p1-p9)
% - Enhanced TikZ/PGFPlots for dark backgrounds
% - qq tcolorbox environment
% - tkz-euclide for geometric constructions

% ===========================================
% DARK MODE PACKAGE
% ===========================================
\usepackage{darkmode}
\enabledarkmode

% ===========================================
% ADDITIONAL PACKAGES
% ===========================================
\usepackage{changepage}
\usepackage{ulem}
\usepackage{colortbl}
\usepackage{tkz-euclide}
\usepackage{capt-of}

% ===========================================
% EXTENDED TIKZ LIBRARIES
% ===========================================
\usetikzlibrary{patterns,3d,backgrounds}
\usepgfplotslibrary{groupplots}

% Upgrade pgfplots compatibility
\pgfplotsset{compat=1.18}

% ===========================================
% OCEAN BLUE GRADIENT (p1-p9)
% ===========================================
\definecolor{p1}{HTML}{caf0f8}
\definecolor{p2}{HTML}{ade8f4}
\definecolor{p3}{HTML}{90e0ef}
\definecolor{p4}{HTML}{48cae4}
\definecolor{p5}{HTML}{00b4d8}
\definecolor{p6}{HTML}{0096c7}
\definecolor{p7}{HTML}{0077b6}
\definecolor{p8}{HTML}{023e8a}
\definecolor{p9}{HTML}{03045e}

% Dark background color
\definecolor{pag}{HTML}{293133}

% ===========================================
% MATERIAL DESIGN COLORS
% ===========================================
% Note: myred, myyellow already defined in main preamble
% These are additional/alternate versions
\definecolor{matred}{HTML}{f44336}
\definecolor{matpink}{HTML}{e81e63}
\definecolor{matpurple}{HTML}{9c27b0}
\definecolor{matdeeppurple}{HTML}{673ab7}
\definecolor{matindigo}{HTML}{3f51b5}
\definecolor{matblue}{HTML}{2196f3}
\definecolor{matlightblue}{HTML}{03a9f4}
\definecolor{matcyan}{HTML}{00bcd4}
\definecolor{matteal}{HTML}{009688}
\definecolor{matgreen}{HTML}{4caf50}
\definecolor{matlightgreen}{HTML}{8bc34a}
\definecolor{matlime}{HTML}{cddc39}
\definecolor{matyellow}{HTML}{ffeb3b}
\definecolor{matamber}{HTML}{ffc107}
\definecolor{matorange}{HTML}{ff9800}
\definecolor{matdeeporange}{HTML}{ff5722}
\definecolor{matbrown}{HTML}{795548}
\definecolor{matgray}{HTML}{9e9e9e}
\definecolor{matbluegray}{HTML}{607d8b}

% ===========================================
% COLORLET FOR TIKZ
% ===========================================
\colorlet{xcol}{blue!60!black}

% ===========================================
% DARK MODE TCOLORBOX
% ===========================================
\newtcolorbox{qq}[2][]{%
    boxrule=0.75pt,
    sharp corners,
    colframe=white,
    colback=pag,
    coltext=white,
    #1,
}

% Dark definition box
\newtcolorbox{darkdfn}[2][]{%
    enhanced,
    boxrule=1pt,
    sharp corners,
    colframe=p5,
    colback=pag,
    coltext=white,
    coltitle=white,
    fonttitle=\bfseries,
    title=#2,
    #1,
}

% Dark theorem box
\newtcolorbox{darkthm}[2][]{%
    enhanced,
    boxrule=1pt,
    sharp corners,
    colframe=matpurple,
    colback=pag,
    coltext=white,
    coltitle=white,
    fonttitle=\bfseries,
    title=#2,
    #1,
}

% ===========================================
% TIKZ 3D FIX
% ===========================================
% Small fix for canvas is xy plane at z
% https://tex.stackexchange.com/a/48776/121799
\makeatletter
\tikzoption{canvas is xy plane at z}[]{%
    \def\tikz@plane@origin{\pgfpointxyz{0}{0}{#1}}%
    \def\tikz@plane@x{\pgfpointxyz{1}{0}{#1}}%
    \def\tikz@plane@y{\pgfpointxyz{0}{1}{#1}}%
    \tikz@canvas@is@plane}
\makeatother

% ===========================================
% CONDITIONAL LINE TIKZSET
% ===========================================
\tikzset{
    conditional line/.style 2 args={
        /utils/exec={\pgfmathsetmacro{\xcoordA}{#1}},
        /utils/exec={\pgfmathsetmacro{\xcoordB}{#2}},
        insert path={
            \ifdim\xcoordA pt>\xcoordB pt
            (\xcoordA,0) -- (\xcoordB,0)
            \fi
        },
    },
}

% ===========================================
% PGFPLOTS DARK MODE DEFAULTS
% ===========================================
\pgfplotsset{
    dark axis/.style={
        axis line style={white},
        tick style={white},
        label style={white},
        grid style={pag!70},
        every axis label/.append style={white},
        every tick label/.append style={white},
    },
}

% ===========================================
% TIKZ EXTERNALIZATION (for fast compilation)
% ===========================================
% This creates .md5 cache files for figures
% Requires: pdflatex -shell-escape
%
% To enable, add AFTER loading this file:
%   \enabletikzexternalize
%
\newcommand{\enabletikzexternalize}{%
  \usetikzlibrary{external}%
  \tikzexternalize[prefix=figures/]%
}

% ===========================================
% CONVENIENT SHORTCUTS FOR DARK PLOTS
% ===========================================
\newcommand{\darkplot}[1]{%
    \begin{tikzpicture}
    \begin{axis}[
        dark axis,
        grid=both,
        major grid style={line width=.2pt,draw=pag!70},
        #1
    ]
}


% ===========================================
% QUICK GEOMETRY MACROS (tkz-euclide)
% ===========================================
% Draw a point: \point{name}{x}{y}
\newcommand{\gpoint}[3]{\tkzDefPoint(#2,#3){#1}}

% Draw line through points: \gline{A}{B}
\newcommand{\gline}[2]{\tkzDrawLine[color=p4](#1,#2)}

% Draw circle: \gcircle{center}{radius}
\newcommand{\gcircle}[2]{\tkzDrawCircle[color=p5](#1,#2)}

% Mark angle: \gangle{A}{B}{C}
\newcommand{\gangle}[3]{\tkzMarkAngle[color=p3,size=0.5](#1,#2,#3)}

% ===========================================
% USAGE EXAMPLE
% ===========================================
% In your document:
%
% \begin{tikzpicture}
%   \begin{axis}[dark axis, ...]
%     \addplot[p4, thick] {sin(deg(x))};
%   \end{axis}
% \end{tikzpicture}
%
% Or with qq box:
% \begin{qq}{My Title}
%   Content here in dark box
% \end{qq}


\course{Studies-in-Algebra-and-Group-Theory}

\begin{document}

\lecture{7}{More maps and spaces}

\subsection{Annihilators}

The annihilator provides a precise duality between subspaces of \(V\) and subspaces of \(V^*\).

\dfn{Annihilator}{For a subspace \(U \subseteq V\), the \textbf{annihilator} of \(U\) is
\[\operatorname{Ann}(U) = U^\circ = \{\ell \in V^* : \ell(u) = 0 \text{ for all } u \in U\}.\]}

\thm{Dimension of the annihilator}{If \(V\) is finite-dimensional and \(U \subseteq V\) is a subspace, then
\[\dim(\operatorname{Ann}(U)) = \dim(V) - \dim(U).\]}

\pf{Proof}{The restriction map \(\rho : V^* \to U^*\) defined by \(\rho(\ell) = \ell|_U\) is a linear map with \(\ker(\rho) = \operatorname{Ann}(U)\). We claim \(\rho\) is surjective: given \(\psi \in U^*\), extend a basis of \(U\) to a basis of \(V\), define \(\ell\) to agree with \(\psi\) on \(U\) and be zero on the complementary basis vectors. Then \(\rho(\ell) = \psi\).

By rank-nullity:
\[\dim(V^*) = \dim(\operatorname{Ann}(U)) + \dim(U^*).\]
Since \(\dim(V^*) = \dim(V)\) and \(\dim(U^*) = \dim(U)\), we get \(\dim(\operatorname{Ann}(U)) = \dim(V) - \dim(U)\).
}

\thm{Annihilator isomorphism}{For a subspace \(U \subseteq V\), there is a natural isomorphism
\[\operatorname{Ann}(U) \cong (V/U)^*.\]}

\pf{Proof}{Define \(\Phi : \operatorname{Ann}(U) \to (V/U)^*\) by \(\Phi(\ell)(v + U) = \ell(v)\). 

\textbf{Well-defined:} If \(v + U = v' + U\), then \(v - v' \in U\), so \(\ell(v) - \ell(v') = \ell(v - v') = 0\) since \(\ell \in \operatorname{Ann}(U)\).

\textbf{Injective:} If \(\Phi(\ell) = 0\), then \(\ell(v) = 0\) for all \(v \in V\), so \(\ell = 0\).

\textbf{Surjective:} Given \(\psi \in (V/U)^*\), define \(\ell = \psi \circ q\) where \(q : V \to V/U\) is the quotient map. Then \(\ell \in \operatorname{Ann}(U)\) and \(\Phi(\ell) = \psi\).
}



\section{Quotient Spaces}

Quotient spaces provide a systematic way to ``collapse'' a subspace to zero, creating a smaller vector space that captures the structure ``perpendicular'' to the subspace. This construction is fundamental in linear algebra and pervades abstract algebra.

\subsection{Construction of the Quotient Space}

\dfn{Cosets and quotient space}{Let \(V\) be a vector space over \(k\) and \(U \subseteq V\) a subspace. For \(v \in V\), the \textbf{coset} of \(v\) modulo \(U\) is
\[v + U = \{v + u : u \in U\}.\]
The \textbf{quotient space} is the set of all cosets:
\[V/U = \{v + U : v \in V\}.\]}

\nt{Two cosets \(v + U\) and \(w + U\) are equal if and only if \(v - w \in U\). Equivalently, the cosets partition \(V\) into equivalence classes under the relation \(v \sim w \Leftrightarrow v - w \in U\).}

\mprop{Vector space structure on the quotient}{The quotient \(V/U\) becomes a vector space over \(k\) under the operations:
\begin{enumerate}
  \item \textbf{Addition:} \((v + U) + (w + U) = (v + w) + U\).
  \item \textbf{Scalar multiplication:} \(\lambda(v + U) = (\lambda v) + U\).
\end{enumerate}
The zero element is \(0 + U = U\), and \(-(v + U) = (-v) + U\).}

\pf{Proof}{The key is \emph{well-definedness}: if \(v + U = v' + U\) and \(w + U = w' + U\), then \(v - v' \in U\) and \(w - w' \in U\), so \((v + w) - (v' + w') = (v - v') + (w - w') \in U\), meaning \((v + w) + U = (v' + w') + U\).

Similarly for scalar multiplication: if \(v + U = v' + U\), then \(v - v' \in U\), so \(\lambda(v - v') = \lambda v - \lambda v' \in U\), giving \(\lambda v + U = \lambda v' + U\).

The vector space axioms follow from those in \(V\).
}

\subsection{The Quotient Map and Exact Sequences}

\dfn{Quotient map}{The \textbf{quotient map} (or \textbf{canonical projection}) is \(q : V \to V/U\) defined by \(q(v) = v + U\).}

\mprop{Properties of the quotient map}{The quotient map \(q : V \to V/U\) satisfies:
\begin{enumerate}
  \item \(q\) is linear.
  \item \(q\) is surjective.
  \item \(\ker(q) = U\).
\end{enumerate}
}


\pf{Proof}{(1) \(q(v + \lambda w) = (v + \lambda w) + U = (v + U) + \lambda(w + U) = q(v) + \lambda q(w)\).

(2) Every coset \(v + U\) is \(q(v)\).

(3) \(v \in \ker(q) \Leftrightarrow v + U = 0 + U \Leftrightarrow v \in U\).
}

\nt{This gives a \textbf{short exact sequence}:
\[0 \longrightarrow U \xlongrightarrow{\iota} V \xlongrightarrow{q} V/U \longrightarrow 0\]
where \(\iota : U \hookrightarrow V\) is the inclusion. ``Exact'' means: \(\Img(\iota) = \ker(q)\), \(\iota\) is injective, and \(q\) is surjective. Short exact sequences are a central organizing principle in algebra.}

\thm{Dimension of the quotient space}{If \(V\) is finite-dimensional and \(U \subseteq V\) is a subspace, then
\[\dim(V/U) = \dim(V) - \dim(U).\]}

\pf{Proof}{Apply rank-nullity to \(q : V \to V/U\):
\[\dim(V) = \dim(\ker(q)) + \dim(\Img(q)) = \dim(U) + \dim(V/U).\]
}

\ex{Concrete example}{Let \(V = \RR^3\) and \(U = \operatorname{Span}\{(1, 0, 0)\}\), the \(x\)-axis. Then \(V/U \cong \RR^2\):

The coset \((a, b, c) + U\) consists of all points \(\{(a + t, b, c) : t \in \RR\}\)---a horizontal line through \((a, b, c)\). Two points have the same coset iff they have the same \(y\) and \(z\) coordinates. The map \((a, b, c) + U \mapsto (b, c)\) is an isomorphism \(V/U \xrightarrow{\sim} \RR^2\).
}

\subsection{Universal Property of the Quotient}

The quotient space is characterized by a universal property that explains its importance.

\thm{Universal property of the quotient}{Let \(\varphi : V \to W\) be a linear map with \(U \subseteq \ker(\varphi)\). Then there exists a unique linear map \(\bar{\varphi} : V/U \to W\) such that \(\varphi = \bar{\varphi} \circ q\):
\[\begin{tikzcd}[ampersand replacement=\&]
V \ar[r, "\varphi"] \ar[d, "q"'] \& W \\
V/U \ar[ur, dashed, "\exists! \bar{\varphi}"']
\end{tikzcd}\]
Moreover, \(\bar{\varphi}\) is defined by \(\bar{\varphi}(v + U) = \varphi(v)\).}

\pf{Proof}{\textbf{Well-defined:} If \(v + U = v' + U\), then \(v - v' \in U \subseteq \ker(\varphi)\), so \(\varphi(v) = \varphi(v')\).

\textbf{Linearity:} \(\bar{\varphi}((v + U) + \lambda(w + U)) = \bar{\varphi}((v + \lambda w) + U) = \varphi(v + \lambda w) = \varphi(v) + \lambda \varphi(w) = \bar{\varphi}(v + U) + \lambda \bar{\varphi}(w + U)\).

\textbf{Uniqueness:} Any \(\psi : V/U \to W\) with \(\varphi = \psi \circ q\) must satisfy \(\psi(v + U) = \psi(q(v)) = \varphi(v) = \bar{\varphi}(v + U)\).
}


\cor{First Isomorphism Theorem}{For any linear map \(\varphi : V \to W\):
\[V/\ker(\varphi) \cong \Img(\varphi).\]
The isomorphism is \(\bar{\varphi} : v + \ker(\varphi) \mapsto \varphi(v)\).}

\pf{Proof}{By the universal property, \(\bar{\varphi} : V/\ker(\varphi) \to W\) is well-defined. Its image is exactly \(\Img(\varphi)\). Injectivity: if \(\bar{\varphi}(v + \ker(\varphi)) = 0\), then \(\varphi(v) = 0\), so \(v \in \ker(\varphi)\), meaning \(v + \ker(\varphi) = 0 + \ker(\varphi)\).
}

\nt{This theorem is remarkably useful. It says that every surjective linear map \(\varphi : V \twoheadrightarrow W\) is ``really'' a quotient map: \(V \to V/\ker(\varphi) \cong W\). Dually, every injective map \(\iota : U \hookrightarrow V\) makes \(U\) isomorphic to a subspace of \(V\).}

\subsection{Correspondence Between Subspaces}

The quotient construction establishes a bijection between certain subspaces.

\thm{Lattice correspondence}{Let \(U \subseteq V\) be a subspace. There is a bijection
\[\{\text{subspaces } W \subseteq V \text{ with } U \subseteq W\} \longleftrightarrow \{\text{subspaces of } V/U\}\]
given by:
\begin{itemize}
  \item Forward: \(W \mapsto W/U = \{w + U : w \in W\}\).
  \item Backward: \(\bar{W} \mapsto q^{-1}(\bar{W}) = \{v \in V : v + U \in \bar{W}\}\).
\end{itemize}
Moreover, this bijection preserves inclusions: \(W_1 \subseteq W_2 \Leftrightarrow W_1/U \subseteq W_2/U\).}

\pf{Proof}{We verify the maps are inverses.

\textbf{Forward then backward:} Starting with \(W \supseteq U\), we get \(W/U\), then \(q^{-1}(W/U) = \{v \in V : v + U \in W/U\}\). But \(v + U \in W/U\) iff \(v + U = w + U\) for some \(w \in W\), iff \(v - w \in U \subseteq W\), iff \(v \in W\). So \(q^{-1}(W/U) = W\).

\textbf{Backward then forward:} Starting with \(\bar{W} \subseteq V/U\), let \(W = q^{-1}(\bar{W})\). Note \(U = q^{-1}(0) \subseteq q^{-1}(\bar{W}) = W\). Then \(W/U = q(W) = q(q^{-1}(\bar{W})) = \bar{W}\) (surjectivity of \(q\) ensures \(q(q^{-1}(S)) = S\) for any \(S \subseteq V/U\)).
}

\nt{Maximal proper subspaces of \(V\) containing \(U\) correspond to maximal proper subspaces of \(V/U\)---i.e., hyperplanes (codimension-1 subspaces). This perspective becomes useful in algebraic geometry and representation theory.}



\section{Linear Operators and Endomorphisms}

When the domain and codomain of a linear map coincide, we gain additional algebraic structure: the ability to compose maps with themselves.

\subsection{The Ring of Endomorphisms}

\dfn{Linear operator / Endomorphism}{A \textbf{linear operator} or \textbf{endomorphism} of a vector space \(V\) is a linear map \(\varphi : V \to V\). We write
\[\End(V) = \Hom(V, V)\]
for the set of all endomorphisms of \(V\).}

\mprop{Ring structure on \(\End(V)\)}{\(\End(V)\) is a \textbf{ring} (with identity) under:
\begin{itemize}
  \item \textbf{Addition:} \((\varphi + \psi)(v) = \varphi(v) + \psi(v)\).
  \item \textbf{Multiplication:} \((\varphi \circ \psi)(v) = \varphi(\psi(v))\) (composition).
  \item \textbf{Identity:} The identity map \(\id_V\).
\end{itemize}
Moreover, \(\End(V)\) is a \(k\)-algebra: multiplication distributes over scalar multiplication.}

\pf{Proof}{Addition makes \(\End(V)\) an abelian group (it's a subgroup of \(\Hom(V, V)\) as a vector space). Composition is associative with identity \(\id_V\). Distributivity: \(\varphi \circ (\psi + \rho) = \varphi \circ \psi + \varphi \circ \rho\) and \((\varphi + \psi) \circ \rho = \varphi \circ \rho + \psi \circ \rho\). These follow from linearity.
}

\nt{The ring \(\End(V)\) is \emph{noncommutative} unless \(\dim(V) \leq 1\). For \(V = k^n\), \(\End(V) \cong M_n(k)\), the ring of \(n \times n\) matrices. The noncommutativity of matrix multiplication reflects the noncommutativity of \(\End(V)\).}

\subsection{Polynomials of Operators}

Since endomorphisms can be composed, we can form ``powers'' and evaluate polynomials.

\dfn{Powers and polynomials of operators}{For \(\varphi \in \End(V)\) and \(n \in \ZZ_{\geq 0}\):
\[\varphi^0 = \id_V, \quad \varphi^n = \underbrace{\varphi \circ \varphi \circ \cdots \circ \varphi}_{n \text{ times}}.\]
For a polynomial \(p(x) = a_0 + a_1 x + \cdots + a_n x^n \in k[x]\), define
\[p(\varphi) = a_0 \id_V + a_1 \varphi + \cdots + a_n \varphi^n \in \End(V).\]}

\mprop{Evaluation homomorphism}{The map \(\ev_\varphi : k[x] \to \End(V)\) defined by \(p(x) \mapsto p(\varphi)\) is a ring homomorphism. That is:
\begin{enumerate}
  \item \((p + q)(\varphi) = p(\varphi) + q(\varphi)\).
  \item \((pq)(\varphi) = p(\varphi) \circ q(\varphi)\).
  \item \(1(\varphi) = \id_V\).
\end{enumerate}
}


\pf{Proof}{(1) and (3) are immediate. For (2), let \(p(x) = \sum a_i x^i\) and \(q(x) = \sum b_j x^j\). Then \((pq)(x) = \sum_{k} c_k x^k\) where \(c_k = \sum_{i+j=k} a_i b_j\). So
\[(pq)(\varphi) = \sum_k c_k \varphi^k = \sum_k \left(\sum_{i+j=k} a_i b_j\right) \varphi^k = \left(\sum_i a_i \varphi^i\right) \circ \left(\sum_j b_j \varphi^j\right) = p(\varphi) \circ q(\varphi).\]
}

\nt{This makes \(\End(V)\) a \(k[x]\)-module: \(p(x) \cdot v = p(\varphi)(v)\). The study of this module structure (particularly the kernel of \(\ev_\varphi\), called the \emph{minimal polynomial}) is central to the theory of canonical forms.}

\section{Invariant Subspaces}

Understanding an operator often begins with finding subspaces on which its behavior is simpler.

\dfn{Invariant subspace}{A subspace \(W \subseteq V\) is \textbf{invariant} under \(\varphi \in \End(V)\) (or \textbf{\(\varphi\)-invariant}) if
\[\varphi(W) \subseteq W,\]
i.e., \(\varphi(w) \in W\) for all \(w \in W\).}

\ex{Basic examples of invariant subspaces}{For any \(\varphi \in \End(V)\):
\begin{itemize}
  \item \(\{0\}\) and \(V\) are always \(\varphi\)-invariant.
  \item \(\ker(\varphi)\) is \(\varphi\)-invariant: if \(\varphi(v) = 0\), then \(\varphi(\varphi(v)) = \varphi(0) = 0\).
  \item \(\Img(\varphi)\) is \(\varphi\)-invariant: \(\varphi(\Img(\varphi)) = \Img(\varphi^2) \subseteq \Img(\varphi)\).
  \item More generally, \(\ker(\varphi^n)\) and \(\Img(\varphi^n)\) are \(\varphi\)-invariant for all \(n \geq 0\).
\end{itemize}
}

\mprop{Restriction to an invariant subspace}{If \(W \subseteq V\) is \(\varphi\)-invariant, the \textbf{restriction} \(\varphi|_W : W \to W\) is a well-defined linear operator on \(W\).}

\subsection{Block Diagonal Structure}

Invariant subspaces lead to simpler matrix representations.

\thm{Direct sum of invariant subspaces}{Suppose \(V = V_1 \oplus V_2\) where both \(V_1\) and \(V_2\) are \(\varphi\)-invariant. Choose bases \(\mathcal{B}_1 = \{e_1, \ldots, e_m\}\) of \(V_1\) and \(\mathcal{B}_2 = \{e_{m+1}, \ldots, e_n\}\) of \(V_2\), forming a basis \(\mathcal{B} = \mathcal{B}_1 \cup \mathcal{B}_2\) of \(V\). Then the matrix of \(\varphi\) with respect to \(\mathcal{B}\) is \textbf{block diagonal}:
\[\mcM(\varphi; \mathcal{B}) = \begin{pmatrix} A_1 & 0 \\ 0 & A_2 \end{pmatrix}\]
where \(A_1 = \mcM(\varphi|_{V_1}; \mathcal{B}_1)\) and \(A_2 = \mcM(\varphi|_{V_2}; \mathcal{B}_2)\).}

\pf{Proof}{For \(i \leq m\): \(\varphi(e_i) \in V_1\) (by \(\varphi\)-invariance), so \(\varphi(e_i)\) is a linear combination of \(e_1, \ldots, e_m\) only. For \(i > m\): \(\varphi(e_i) \in V_2\), so \(\varphi(e_i)\) is a linear combination of \(e_{m+1}, \ldots, e_n\) only. This gives the block structure.
}


\nt{Conversely, if a matrix is block diagonal with blocks \(A_1\) and \(A_2\), then the spans of the corresponding basis vectors are invariant subspaces. The goal of much of spectral theory is to decompose \(V\) into a direct sum of invariant subspaces on which \(\varphi\) acts simply.}

\cor{Block triangular from one invariant subspace}{If \(V_1\) is \(\varphi\)-invariant (but we don't assume \(V_1\) has an invariant complement), then extending a basis of \(V_1\) to a basis of \(V\) gives a \textbf{block upper-triangular} matrix:
\[\mcM(\varphi) = \begin{pmatrix} A_1 & * \\ 0 & A_2 \end{pmatrix}.\]
The upper-right block \(*\) may be nonzero; it measures the ``interaction'' between \(V_1\) and \(V/V_1\).}

\section{Eigenvalues and Eigenvectors}

The simplest invariant subspaces are one-dimensional. These give rise to the fundamental notion of eigenvalues.

\subsection{Definitions and Basic Properties}

\dfn{Eigenvector and eigenvalue}{Let \(\varphi \in \End(V)\). A nonzero vector \(v \in V\) is an \textbf{eigenvector} of \(\varphi\) if
\[\varphi(v) = \lambda v\]
for some scalar \(\lambda \in k\). The scalar \(\lambda\) is called the \textbf{eigenvalue} corresponding to \(v\).}

\nt{The condition \(\varphi(v) = \lambda v\) says that \(\varphi\) acts on the line spanned by \(v\) as scaling by \(\lambda\). Equivalently, \(\operatorname{Span}\{v\}\) is a one-dimensional \(\varphi\)-invariant subspace. Conversely, any one-dimensional invariant subspace \(U = \operatorname{Span}\{v\}\) consists of eigenvectors (with the same eigenvalue) plus zero.}

\mprop{Characterization of eigenvalues}{For \(\varphi \in \End(V)\) and \(\lambda \in k\), the following are equivalent:
\begin{enumerate}
  \item \(\lambda\) is an eigenvalue of \(\varphi\).
  \item \(\varphi - \lambda \id_V\) is not injective.
  \item \(\ker(\varphi - \lambda \id_V) \neq \{0\}\).
  \item (If \(V\) is finite-dimensional) \(\varphi - \lambda \id_V\) is not surjective.
  \item (If \(V\) is finite-dimensional) \(\det(\varphi - \lambda \id_V) = 0\).
\end{enumerate}
}

\pf{Proof}{(1) \(\Leftrightarrow\) (2) \(\Leftrightarrow\) (3): \(\lambda\) is an eigenvalue iff there exists nonzero \(v\) with \(\varphi(v) = \lambda v\), iff \((\varphi - \lambda \id)(v) = 0\) for some nonzero \(v\), iff \(\ker(\varphi - \lambda \id) \neq \{0\}\).

(3) \(\Leftrightarrow\) (4) \(\Leftrightarrow\) (5): In finite dimensions, a linear operator is injective iff surjective iff invertible iff has nonzero determinant. So \(\ker(\varphi - \lambda \id) \neq \{0\}\) iff \(\det(\varphi - \lambda \id) = 0\).
}

\dfn{Eigenspace}{For an eigenvalue \(\lambda\) of \(\varphi\), the \textbf{eigenspace} is
\[E(\lambda, \varphi) = \ker(\varphi - \lambda \id_V) = \{v \in V : \varphi(v) = \lambda v\}.\]
This is a subspace of \(V\) (it includes 0, though 0 is not an eigenvector).}


\thm{Eigenvectors with distinct eigenvalues are linearly independent}{Let \(v_1, \ldots, v_m\) be eigenvectors of \(\varphi\) with distinct eigenvalues \(\lambda_1, \ldots, \lambda_m\). Then \(v_1, \ldots, v_m\) are linearly independent.}

\pf{Proof}{Induction on \(m\). Base case \(m = 1\): a single eigenvector is nonzero, hence linearly independent.

Suppose the result holds for \(m - 1\) eigenvectors. Assume \(\sum_{i=1}^m c_i v_i = 0\). Apply \(\varphi\):
\[\sum_{i=1}^m c_i \lambda_i v_i = 0.\]
Multiply the original equation by \(\lambda_m\) and subtract:
\[\sum_{i=1}^{m-1} c_i (\lambda_i - \lambda_m) v_i = 0.\]
By the induction hypothesis (since \(v_1, \ldots, v_{m-1}\) are eigenvectors with distinct eigenvalues \(\lambda_1, \ldots, \lambda_{m-1}\)), we have \(c_i(\lambda_i - \lambda_m) = 0\) for \(i < m\). Since \(\lambda_i \neq \lambda_m\), we get \(c_i = 0\). Substituting back: \(c_m v_m = 0\), so \(c_m = 0\) (since \(v_m \neq 0\)).
}

\cor{Bound on number of eigenvalues}{A linear operator on an \(n\)-dimensional space has at most \(n\) distinct eigenvalues.}

\pf{Proof}{If \(\lambda_1, \ldots, \lambda_m\) are distinct eigenvalues with eigenvectors \(v_1, \ldots, v_m\), then \(v_1, \ldots, v_m\) are linearly independent. In an \(n\)-dimensional space, any linearly independent set has at most \(n\) elements.
}

\cor{Sum of eigenspaces is direct}{If \(\lambda_1, \ldots, \lambda_m\) are distinct eigenvalues of \(\varphi\), then
\[E(\lambda_1, \varphi) + \cdots + E(\lambda_m, \varphi) = E(\lambda_1, \varphi) \oplus \cdots \oplus E(\lambda_m, \varphi).\]
is a direct sum. Furthermore, \(\sum_i \dim E(\lambda_i, \varphi) \leq \dim V\).}

\subsection{Diagonalizable Operators}

The best-case scenario: \(V\) has a basis of eigenvectors.

\dfn{Diagonalizable operator}{An operator \(\varphi \in \End(V)\) is \textbf{diagonalizable} if there exists a basis of \(V\) consisting entirely of eigenvectors of \(\varphi\). Equivalently, \(\varphi\) has a diagonal matrix with respect to some basis.}

\nt{If \(\{v_1, \ldots, v_n\}\) is a basis of eigenvectors with \(\varphi(v_i) = \lambda_i v_i\), then the matrix of \(\varphi\) is
\[\mcM(\varphi) = \begin{pmatrix} \lambda_1 & & 0 \\ & \ddots & \\ 0 & & \lambda_n \end{pmatrix}.\]
The eigenvalues appear on the diagonal (possibly with repetition).}

\thm{Characterizations of diagonalizability}{For \(\varphi \in \End(V)\) with \(\dim(V) = n\) and distinct eigenvalues \(\lambda_1, \ldots, \lambda_m\), the following are equivalent:
\begin{enumerate}
  \item \(\varphi\) is diagonalizable.
  \item \(V\) has a basis of eigenvectors of \(\varphi\).
  \item \(V = E(\lambda_1, \varphi) \oplus \cdots \oplus E(\lambda_m, \varphi)\).
  \item \(\dim E(\lambda_1, \varphi) + \cdots + \dim E(\lambda_m, \varphi) = n\).
\end{enumerate}
}


\pf{Proof}{(1) \(\Leftrightarrow\) (2): By definition.

(2) \(\Rightarrow\) (3): If \(V\) has a basis of eigenvectors, partition it by eigenvalue. Each part is a basis of the corresponding eigenspace, and together they span \(V\).

(3) \(\Rightarrow\) (4): Take dimensions: \(\dim V = \dim(E_1 \oplus \cdots \oplus E_m) = \sum \dim E_i\).

(4) \(\Rightarrow\) (2): The sum \(E_1 + \cdots + E_m\) is direct (earlier corollary), and has dimension \(\sum \dim E_i = n = \dim V\). So \(E_1 \oplus \cdots \oplus E_m = V\). Combining bases of each \(E_i\) gives a basis of \(V\) consisting of eigenvectors.
}

\cor{Enough distinct eigenvalues implies diagonalizability}{If \(\varphi \in \End(V)\) has \(\dim(V)\) distinct eigenvalues, then \(\varphi\) is diagonalizable.}

\pf{Proof}{Each eigenspace has dimension at least 1 (it contains an eigenvector). With \(n = \dim V\) distinct eigenvalues, \(\sum \dim E(\lambda_i) \geq n\). But this sum is also \(\leq n\), so equality holds, and criterion (4) above is satisfied.
}

\subsection{Examples: Diagonalizable and Non-Diagonalizable}

\ex{Diagonal matrices}{The identity matrix \(I_n\) has eigenvalue \(1\) with eigenspace all of \(k^n\). More generally, a diagonal matrix \(\diag(\lambda_1, \ldots, \lambda_n)\) has the standard basis vectors as eigenvectors with eigenvalues \(\lambda_1, \ldots, \lambda_n\). Such a matrix is (trivially) diagonalizable---it's already diagonal!}

\ex{A non-diagonalizable matrix}{Consider \(\varphi : \RR^2 \to \RR^2\) with matrix
\[A = \begin{pmatrix} 1 & 1 \\ 0 & 1 \end{pmatrix}.\]
The eigenvalues satisfy \(\det(A - \lambda I) = (1 - \lambda)^2 = 0\), so \(\lambda = 1\) is the only eigenvalue.

The eigenspace: \(E(1, \varphi) = \ker(A - I) = \ker\begin{pmatrix} 0 & 1 \\ 0 & 0 \end{pmatrix} = \operatorname{Span}\{e_1\}\).

This has dimension 1, but \(\dim(\RR^2) = 2\). So \(\varphi\) is \textbf{not diagonalizable}.

Geometrically, \(\varphi\) is a ``shear'': it fixes the \(x\)-axis pointwise and slides other points parallel to the \(x\)-axis. The only eigenvector direction is along the \(x\)-axis.}

\ex{No eigenvectors over \(\RR\)}{Consider rotation by $90^\circ$ in \(\RR^2\):
\[R = \begin{pmatrix} 0 & -1 \\ 1 & 0 \end{pmatrix}.\]
The characteristic polynomial is \(\det(R - \lambda I) = \lambda^2 + 1\), which has no real roots. So \(R\) has \emph{no eigenvalues} over \(\RR\), hence no eigenvectors. Over \(\CC\), the eigenvalues are \(\pm i\) with eigenvectors \((1, \mp i)^T\), and \(R\) is diagonalizable as a complex matrix. The eigenfairy doesn't visit \( \mathbb{R} \) (sad).

This illustrates that diagonalizability depends on the scalar field!}

\nt{The rotation example shows that not every operator is diagonalizable. Even over \(\CC\), where every polynomial has roots, an operator may fail to be diagonalizable (like the shear above). Understanding non-diagonalizable operators requires the theory of Jordan normal form, which we'll develop later.}


\qs{Exercises on quotient spaces, operators, and eigenvalues}{
\begin{enumerate}
  \item Let \(V = \RR^4\) and \(U = \{(x, y, z, w) : x + y = 0, z = w\}\). Find \(\dim(U)\), \(\dim(V/U)\), and describe \(V/U\) explicitly.
  
  \item Prove that if \(\varphi, \psi \in \End(V)\) commute (\(\varphi \circ \psi = \psi \circ \varphi\)), then every eigenspace of \(\varphi\) is \(\psi\)-invariant.
  
  \item Let \(\varphi \in \End(V)\) be diagonalizable with eigenvalues \(\lambda_1, \ldots, \lambda_m\). Prove that \(p(\varphi) = 0\) where \(p(x) = \prod_{i=1}^m (x - \lambda_i)\).
  
  \item Show that the quotient map \(q : V \to V/U\) induces a bijection between \(\{\text{linear maps } V/U \to W\}\) and \(\{\text{linear maps } V \to W \text{ vanishing on } U\}\).
  
  \item Let \(A = \begin{pmatrix} 4 & 2 \\ 1 & 3 \end{pmatrix}\). Find the eigenvalues and eigenspaces of \(A\). Is \(A\) diagonalizable? If so, find a matrix \(P\) such that \(P^{-1}AP\) is diagonal.
\end{enumerate}
}

\begin{solution}
\textbf{(1)} The constraints \(x + y = 0\) and \(z = w\) give 2 independent linear equations on \(\RR^4\), so \(\dim(U) = 4 - 2 = 2\). A basis: \(\{(-1, 1, 0, 0), (0, 0, 1, 1)\}\).

By the dimension formula: \(\dim(V/U) = \dim(V) - \dim(U) = 4 - 2 = 2\).

The map \(\pi : \RR^4 \to \RR^2\) defined by \(\pi(x, y, z, w) = (x + y, z - w)\) has \(\ker(\pi) = U\). By the first isomorphism theorem, \(V/U \cong \RR^2\). Explicitly, the coset \((x, y, z, w) + U\) is determined by \((x + y, z - w)\).

\textbf{(2)} Let \(v\) be an eigenvector of \(\varphi\) with eigenvalue \(\lambda\): \(\varphi(v) = \lambda v\). We need \(\psi(v) \in E(\lambda, \varphi)\), i.e., \(\varphi(\psi(v)) = \lambda \psi(v)\).

Compute: \(\varphi(\psi(v)) = \psi(\varphi(v)) = \psi(\lambda v) = \lambda \psi(v)\). $\checkmark$

(We used commutativity in the first equality.)

\textbf{(3)} Let \(p(x) = \prod_{i=1}^m (x - \lambda_i)\). Since \(\varphi\) is diagonalizable, \(V = \bigoplus_{i=1}^m E(\lambda_i, \varphi)\). For any eigenvector \(v \in E(\lambda_j, \varphi)\):
\[(\varphi - \lambda_j \id)(v) = 0.\]
So \(p(\varphi)(v) = (\varphi - \lambda_1) \cdots (\varphi - \lambda_m)(v) = 0\) because one factor gives zero.

Since this holds for all eigenvectors and they span \(V\), we have \(p(\varphi) = 0\).

\textbf{(4)} Define \(\Phi : \Hom(V/U, W) \to \{T \in \Hom(V, W) : U \subseteq \ker(T)\}\) by \(\Phi(\bar{T}) = \bar{T} \circ q\).

\emph{Injective:} If \(\bar{T} \circ q = 0\), then \(\bar{T}(v + U) = \bar{T}(q(v)) = 0\) for all \(v\), so \(\bar{T} = 0\).

\emph{Surjective:} Given \(T : V \to W\) with \(U \subseteq \ker(T)\), define \(\bar{T} : V/U \to W\) by \(\bar{T}(v + U) = T(v)\). Well-defined because \(v - v' \in U \Rightarrow T(v) = T(v')\). Then \(\Phi(\bar{T}) = \bar{T} \circ q = T\).

\textbf{(5)} Characteristic polynomial: \(\det(A - \lambda I) = (4 - \lambda)(3 - \lambda) - 2 = \lambda^2 - 7\lambda + 10 = (\lambda - 5)(\lambda - 2)\).

Eigenvalues: \(\lambda_1 = 5\), \(\lambda_2 = 2\).

Eigenspaces:
\begin{itemize}
\item \(\lambda = 5\): \((A - 5I)v = 0\) gives \(\begin{pmatrix} -1 & 2 \\ 1 & -2 \end{pmatrix}\begin{pmatrix} x \\ y \end{pmatrix} = 0\), so \(x = 2y\). Eigenspace: \(\operatorname{Span}\{(2, 1)^T\}\).
\item \(\lambda = 2\): \((A - 2I)v = 0\) gives \(\begin{pmatrix} 2 & 2 \\ 1 & 1 \end{pmatrix}\begin{pmatrix} x \\ y \end{pmatrix} = 0\), so \(x = -y\). Eigenspace: \(\operatorname{Span}\{(1, -1)^T\}\).
\end{itemize}

Since we have 2 linearly independent eigenvectors in \(\RR^2\), \(A\) is diagonalizable.

Let \(P = \begin{pmatrix} 2 & 1 \\ 1 & -1 \end{pmatrix}\) (columns are eigenvectors). Then \(P^{-1}AP = \begin{pmatrix} 5 & 0 \\ 0 & 2 \end{pmatrix}\).
\end{solution}


\section{Let's get general up in here!}

\subsection{Generalized Kernals (Null Spaces)}
We want to look at general eigenvectors: things that are not \textit{actually} eigenvectors but "kinda want to be" in a generalized sense. What in tarnation does that mean!? Simplest Eigenspace is just \( 0 \), so let's focus on that to get an understanding:

\dfn{Generalized kernal}{The \emph{generalized kernal} of \( \varphi \) is \( \text{gker}(\varphi ) = \{v \in  V \mid \exists m>0 \text{ s.t. } \varphi ^m(v)=0\}  \). These are all the vecrtors that are evnetually sent to \( 0 \) by repeatedly applying \( \varphi  \).}

\nt{We note that \( \text{gker}(\varphi ) \) is not standard notation. Auroux likes it though, so that is what we are going with! \textit{If you meet a stranger on the street and bring up gker be sure to tell them what you mean, lest they be confused}.}

The difference between the kernal and generalized kernal can be summed up by this: you are in the kernal if you go to \( 0 \) in one step; if you apply itterative steps---\( \varphi ^m \)---and end up at \( 0 \) you are in the generalized kernal. If I look at the kernal of \( \varphi  \) I notice that it is in the kernal of \( \varphi ^2 \); then too this holds for more! They may or may not get larger but they are subsets:\[
  \Ker(\varphi ) \subset \Ker (\varphi ^2) \subset \Ker \varphi^3  \subset \ldots
\]
\mprop{}{If \( \Ker (\varphi ^m) = \Ker (\varphi^{m+1}) \) then there is no strictly larger kernal.}
\pf{Proof}{\( \varphi ^{m+2}(v)=\varphi ^{m+1}(\varphi (v))=0 \iff \varphi (v)\in \Ker \varphi ^m \iff \varphi ^m(\varphi (v))=\varphi ^{m+1}(v)=0 \). So \( \Ker (\varphi ^{m+1})=\Ker (\varphi ^{m+2}).\) Since the sequence stops increasing after at most \( n= \dim V \) steps, \( \text{gker}(\varphi )=\Ker \varphi ^n \).}

\ex{}{\( \varphi : k^2 \to k^2  \) with \( e_{1} \mapsto 0, \, e_{2} \mapsto e_{1} \) and \(
      \begin{pmatrix} 0 & 1 \\ 0 & 0 \end{pmatrix} 
  ,\) then \( \Ker (\varphi )=k \cdot e_{1}, \) but \( \Ker (\varphi ^2) =\text{gker}(\varphi )=k^2\).
}

\nt{The generalized kernal is often referred to as the \emph{null space} (this is the terminology that Axler uses, for example). I may use this interchangeably. }

Since it is not the case that \( V = \text{gker}(\varphi )\oplus \text{Img}(\varphi ) \) for every \( \varphi \in \Hom(V, V) \), we can use the following proposition as a substitute:

\mprop{Generalized kernel-image decomposition}{Suppose \( \varphi \in \End(V) \) and let \(m\) be the smallest positive integer such that \(\Ker(\varphi^m) = \Ker(\varphi^{m+1})\). Then
\[
    V = \Ker(\varphi^m) \oplus \Img(\varphi^m).
\]
In particular, \( V = \text{gker}(\varphi) \oplus \Img(\varphi^{\dim V}) \).}

\pf{Proof}{We verify the two conditions for a direct sum decomposition.

\textbf{Intersection is trivial:} Suppose \(v \in \Ker(\varphi^m) \cap \Img(\varphi^m)\). Write \(v = \varphi^m(u)\) for some \(u \in V\). Then \(\varphi^m(v) = \varphi^{2m}(u) = 0\), so \(u \in \Ker(\varphi^{2m})\). 

We claim \(\Ker(\varphi^{2m}) = \Ker(\varphi^m)\). By the stabilization property, once the kernel sequence stabilizes at step \(m\), it stays constant: \(\Ker(\varphi^m) = \Ker(\varphi^{m+1}) = \cdots = \Ker(\varphi^{2m})\). Thus \(u \in \Ker(\varphi^m)\), which means \(v = \varphi^m(u) = 0\).

\textbf{Sum is all of \(V\):} By rank-nullity applied to \(\varphi^m\):
\[\dim V = \dim \Ker(\varphi^m) + \dim \Img(\varphi^m).\]
Since we just showed \(\Ker(\varphi^m) \cap \Img(\varphi^m) = \{0\}\), the sum \(\Ker(\varphi^m) + \Img(\varphi^m)\) is direct, and its dimension equals \(\dim \Ker(\varphi^m) + \dim \Img(\varphi^m) = \dim V\). Hence it equals \(V\).
}

\nt{This decomposition is the key structural result for understanding non-diagonalizable operators. The generalized kernel \(\text{gker}(\varphi)\) captures the ``nilpotent part'' of \(\varphi\), while \(\Img(\varphi^m)\) is where \(\varphi\) acts invertibly. On \(\text{gker}(\varphi)\), repeated application of \(\varphi\) eventually kills everything; on \(\Img(\varphi^m)\), the restriction of \(\varphi\) is an isomorphism.}

\cor{Restriction to image is an isomorphism}{With notation as above, the restriction \(\varphi|_{\Img(\varphi^m)} : \Img(\varphi^m) \to \Img(\varphi^m)\) is an isomorphism.}

\pf{Proof}{First, \(\Img(\varphi^m)\) is \(\varphi\)-invariant since \(\varphi(\Img(\varphi^m)) = \Img(\varphi^{m+1}) \subseteq \Img(\varphi^m)\). For injectivity: if \(v \in \Img(\varphi^m)\) and \(\varphi(v) = 0\), then \(v \in \Ker(\varphi) \subseteq \Ker(\varphi^m)\). But \(\Ker(\varphi^m) \cap \Img(\varphi^m) = \{0\}\), so \(v = 0\). Since \(\varphi|_{\Img(\varphi^m)}\) is injective on a finite-dimensional space, it's also surjective onto its image, which equals \(\Img(\varphi^m)\) by \(\varphi\)-invariance.
}


\subsection{Generalized Eigenspaces}

Now we extend the idea of generalized kernels to arbitrary eigenvalues, not just \(\lambda = 0\).

\dfn{Generalized eigenspace}{Let \(\varphi \in \End(V)\) and let \(\lambda\) be an eigenvalue of \(\varphi\). The \textbf{generalized eigenspace} of \(\varphi\) corresponding to \(\lambda\) is
\[V_\lambda = \Ker(\varphi - \lambda \id)^{\dim V} = \{v \in V : (\varphi - \lambda \id)^n(v) = 0 \text{ for some } n \geq 1\}.\]
Equivalently, \(V_\lambda = \text{gker}(\varphi - \lambda \id)\).}

\nt{The ordinary eigenspace \(E(\lambda, \varphi) = \Ker(\varphi - \lambda \id)\) is always contained in \(V_\lambda\), but \(V_\lambda\) may be strictly larger. The elements of \(V_\lambda \setminus E(\lambda, \varphi)\) are called \textbf{generalized eigenvectors}---they are not true eigenvectors, but they are ``eventually killed'' by \(\varphi - \lambda \id\).}

\ex{Generalized eigenspace of the shear}{Recall the shear \(\varphi : \RR^2 \to \RR^2\) with matrix \(A = \begin{pmatrix} 1 & 1 \\ 0 & 1 \end{pmatrix}\). The only eigenvalue is \(\lambda = 1\), and the eigenspace is \(E(1, \varphi) = \Ker(A - I) = \Ker\begin{pmatrix} 0 & 1 \\ 0 & 0 \end{pmatrix} = \operatorname{Span}\{e_1\}\), which is one-dimensional.

However, \((A - I)^2 = \begin{pmatrix} 0 & 1 \\ 0 & 0 \end{pmatrix}^2 = \begin{pmatrix} 0 & 0 \\ 0 & 0 \end{pmatrix} = 0\).

So the generalized eigenspace is \(V_1 = \Ker(A - I)^2 = \RR^2\). The vector \(e_2\) is a generalized eigenvector: \((A - I)e_2 = e_1 \neq 0\), but \((A - I)^2 e_2 = 0\).
}

\mprop{Basic properties of generalized eigenspaces}{Let \(\varphi \in \End(V)\) with eigenvalue \(\lambda\). Then:
\begin{enumerate}
    \item \(E(\lambda, \varphi) \subseteq V_\lambda\).
    \item \(V_\lambda\) is \(\varphi\)-invariant.
    \item On \(V_\lambda\), the operator \(\varphi - \lambda \id\) is nilpotent.
    \item If \(\mu \neq \lambda\) are distinct eigenvalues, then \(V_\lambda \cap V_\mu = \{0\}\).
\end{enumerate}
}

\pf{Proof}{(1) is immediate: if \((\varphi - \lambda \id)(v) = 0\), then certainly \((\varphi - \lambda \id)^n(v) = 0\) for all \(n \geq 1\).

(2) Let \(v \in V_\lambda\), so \((\varphi - \lambda \id)^n(v) = 0\) for some \(n\). Since \(\varphi\) commutes with \(\varphi - \lambda \id\), we have
\[(\varphi - \lambda \id)^n(\varphi(v)) = \varphi((\varphi - \lambda \id)^n(v)) = \varphi(0) = 0.\]
Thus \(\varphi(v) \in V_\lambda\).

(3) By definition, \(V_\lambda = \Ker(\varphi - \lambda \id)^m\) for some \(m \leq \dim V\). So \((\varphi - \lambda \id)^m\) restricted to \(V_\lambda\) is the zero map.

(4) Suppose \(v \in V_\lambda \cap V_\mu\) with \(\lambda \neq \mu\). Then there exist \(j, k \geq 1\) such that \((\varphi - \lambda \id)^j(v) = 0\) and \((\varphi - \mu \id)^k(v) = 0\).

The polynomials \(p(x) = (x - \lambda)^j\) and \(q(x) = (x - \mu)^k\) are coprime (they share no common roots). By Bézout's identity, there exist polynomials \(a(x), b(x) \in k[x]\) such that \(a(x)p(x) + b(x)q(x) = 1\).

Evaluating at \(\varphi\): \(a(\varphi)p(\varphi) + b(\varphi)q(\varphi) = \id\). Applying to \(v\):
\[v = a(\varphi)p(\varphi)(v) + b(\varphi)q(\varphi)(v) = a(\varphi)(0) + b(\varphi)(0) = 0.\]
}

\nt{Part (4) shows that generalized eigenspaces for distinct eigenvalues intersect trivially, just like ordinary eigenspaces. The proof uses the coprimality of \((x - \lambda)^j\) and \((x - \mu)^k\)---this is a technique that will reappear in the generalized eigenspace decomposition.}


\subsection{The Generalized Eigenspace Decomposition}

The following theorem is the main structural result for operators on finite-dimensional vector spaces over algebraically closed fields.

\thm{Generalized eigenspace decomposition}{Let \(V\) be a finite-dimensional vector space over an algebraically closed field \(k\), and let \(\varphi \in \End(V)\). Let \(\lambda_1, \ldots, \lambda_m\) be the distinct eigenvalues of \(\varphi\). Then
\[V = V_{\lambda_1} \oplus V_{\lambda_2} \oplus \cdots \oplus V_{\lambda_m}.\]
Moreover, \(\sum_{i=1}^m \dim V_{\lambda_i} = \dim V\).}

\pf{Proof}{We proceed in two steps.

\textbf{Step 1: The sum is direct.} We already showed that \(V_{\lambda_i} \cap V_{\lambda_j} = \{0\}\) for \(i \neq j\). More generally, suppose \(v_1 + \cdots + v_m = 0\) where \(v_i \in V_{\lambda_i}\). We show each \(v_i = 0\).

For each \(i\), let \(n_i\) be such that \((\varphi - \lambda_i \id)^{n_i}(v_i) = 0\). Define \(p_i(x) = (x - \lambda_i)^{n_i}\) and 
\[q_i(x) = \prod_{j \neq i} (x - \lambda_j)^{n_j}.\]
Since \(\lambda_i\) is not a root of \(q_i(x)\), the polynomials \(p_i(x)\) and \(q_i(x)\) are coprime. By Bézout, there exist \(a_i(x), b_i(x)\) with \(a_i(x)p_i(x) + b_i(x)q_i(x) = 1\).

Now, \(q_i(\varphi)(v_j) = 0\) for all \(j \neq i\) (since \(p_j\) divides \(q_i\)). Applying \(q_i(\varphi)\) to \(v_1 + \cdots + v_m = 0\):
\[0 = q_i(\varphi)(v_1 + \cdots + v_m) = q_i(\varphi)(v_i).\]
Thus \(v_i \in \Ker(q_i(\varphi)) \cap \Ker(p_i(\varphi)) = V_{\lambda_i} \cap \Ker(q_i(\varphi))\).

But from \(a_i(\varphi)p_i(\varphi)(v_i) + b_i(\varphi)q_i(\varphi)(v_i) = v_i\) and both terms being zero, we get \(v_i = 0\).

\textbf{Step 2: The sum equals \(V\).} Define \(p(x) = \prod_{i=1}^m (x - \lambda_i)^{n_i}\) where \(n_i = \dim V_{\lambda_i}\) (we'll justify this choice shortly). We claim \(p(\varphi) = 0\), i.e., \(p\) annihilates \(\varphi\).

Consider the polynomial \(q_i(x) = p(x)/(x - \lambda_i)^{n_i}\). These polynomials are pairwise coprime (each \(q_i\) excludes exactly the factor \((x - \lambda_i)^{n_i}\)). Their gcd is 1, so there exist \(r_i(x)\) with \(\sum_i r_i(x)q_i(x) = 1\).

For any \(v \in V\), write \(v = \sum_i r_i(\varphi)q_i(\varphi)(v)\). We claim \(r_i(\varphi)q_i(\varphi)(v) \in V_{\lambda_i}\). Indeed:
\[p_i(\varphi)(r_i(\varphi)q_i(\varphi)(v)) = r_i(\varphi)(p_i(\varphi)q_i(\varphi))(v) = r_i(\varphi)p(\varphi)(v).\]
Once we establish that \(p(\varphi) = 0\) (which follows from Cayley-Hamilton, proved later, or can be verified by dimension counting), this shows \(v\) is a sum of elements from the \(V_{\lambda_i}\).

\textbf{Alternative dimension argument:} Since \(k\) is algebraically closed, \(\varphi\) has at least one eigenvalue. Apply the decomposition \(V = \text{gker}(\varphi - \lambda_1 \id) \oplus \Img((\varphi - \lambda_1 \id)^{n_1})\) where \(n_1\) is the stabilization index. The first summand is \(V_{\lambda_1}\). On the second summand, \(\varphi - \lambda_1 \id\) is invertible, so \(\lambda_1\) is not an eigenvalue of \(\varphi\) restricted there. Proceed by induction on \(\dim V\).
}

\cor{Dimension of generalized eigenspaces}{With notation as above, if \(\chi_\varphi(x) = \prod_{i=1}^m (x - \lambda_i)^{n_i}\) is the characteristic polynomial, then \(\dim V_{\lambda_i} = n_i\).}

\nt{The exponent \(n_i\) in the characteristic polynomial is called the \textbf{algebraic multiplicity} of \(\lambda_i\). The dimension of the ordinary eigenspace \(\dim E(\lambda_i, \varphi)\) is called the \textbf{geometric multiplicity}. We always have:
\[1 \leq \dim E(\lambda_i, \varphi) \leq \dim V_{\lambda_i} = n_i.\]
The operator \(\varphi\) is diagonalizable if and only if geometric multiplicity equals algebraic multiplicity for every eigenvalue.}

\ex{Decomposition for a \(3 \times 3\) matrix}{Consider \(\varphi : \CC^3 \to \CC^3\) with matrix
\[A = \begin{pmatrix} 2 & 1 & 0 \\ 0 & 2 & 0 \\ 0 & 0 & 3 \end{pmatrix}.\]
The characteristic polynomial is \(\chi_A(x) = (x-2)^2(x-3)\), so \(\lambda_1 = 2\) has algebraic multiplicity 2 and \(\lambda_2 = 3\) has algebraic multiplicity 1.

For \(\lambda_1 = 2\): 
\[A - 2I = \begin{pmatrix} 0 & 1 & 0 \\ 0 & 0 & 0 \\ 0 & 0 & 1 \end{pmatrix}, \quad (A-2I)^2 = \begin{pmatrix} 0 & 0 & 0 \\ 0 & 0 & 0 \\ 0 & 0 & 1 \end{pmatrix}.\]
So \(V_2 = \Ker(A-2I)^2 = \operatorname{Span}\{e_1, e_2\}\), which has dimension 2.

For \(\lambda_2 = 3\): \(V_3 = \Ker(A - 3I) = \operatorname{Span}\{e_3\}\), dimension 1.

Indeed, \(\CC^3 = V_2 \oplus V_3\).
}


\section{Nilpotent Operators}

Nilpotent operators are the building blocks for understanding non-diagonalizable maps. On each generalized eigenspace \(V_\lambda\), the operator \(\varphi - \lambda \id\) is nilpotent, so understanding nilpotent operators is the key to understanding all operators.

\subsection{Definition and Basic Properties}

\dfn{Nilpotent operator}{An operator \(\varphi \in \End(V)\) is \textbf{nilpotent} if \(\varphi^m = 0\) for some positive integer \(m\). The smallest such \(m\) is called the \textbf{index of nilpotency} (or \textbf{nilpotence index}) of \(\varphi\).}

\ex{Basic examples of nilpotent operators}{
\begin{enumerate}
    \item The zero operator is nilpotent with index 1.
    \item The shift operator on \(k^n\) with matrix \(\begin{pmatrix} 0 & 1 & 0 & \cdots & 0 \\ 0 & 0 & 1 & \cdots & 0 \\ \vdots & & \ddots & \ddots & \vdots \\ 0 & 0 & \cdots & 0 & 1 \\ 0 & 0 & \cdots & 0 & 0 \end{pmatrix}\) is nilpotent with index \(n\).
    \item Any strictly upper-triangular matrix (zeros on and below the diagonal) is nilpotent.
\end{enumerate}
}

\mprop{Characterizations of nilpotent operators}{For \(\varphi \in \End(V)\) with \(\dim V = n\), the following are equivalent:
\begin{enumerate}
    \item \(\varphi\) is nilpotent.
    \item \(\varphi^n = 0\).
    \item The only eigenvalue of \(\varphi\) is \(0\) (if \(k\) is algebraically closed).
    \item \(V = \text{gker}(\varphi)\).
\end{enumerate}
}

\pf{Proof}{(1) \(\Rightarrow\) (2): If \(\varphi^m = 0\), then the chain \(\{0\} \subseteq \Ker(\varphi) \subseteq \Ker(\varphi^2) \subseteq \cdots\) stabilizes by step \(m\), hence by step \(n\) (since dimensions are bounded by \(n\)). So \(\Ker(\varphi^n) = \text{gker}(\varphi)\). But \(\varphi^m = 0\) means \(\Ker(\varphi^m) = V\), so \(\text{gker}(\varphi) = V\), which means \(\varphi^n = 0\).

(2) \(\Rightarrow\) (3): Suppose \(\lambda\) is an eigenvalue with eigenvector \(v \neq 0\). Then \(\varphi(v) = \lambda v\), so \(\varphi^n(v) = \lambda^n v\). If \(\varphi^n = 0\), then \(\lambda^n v = 0\), so \(\lambda^n = 0\), hence \(\lambda = 0\).

(3) \(\Rightarrow\) (4): Over an algebraically closed field, by the generalized eigenspace decomposition, \(V = \bigoplus_{\lambda} V_\lambda\). If the only eigenvalue is \(0\), then \(V = V_0 = \text{gker}(\varphi)\).

(4) \(\Rightarrow\) (1): If \(V = \text{gker}(\varphi) = \Ker(\varphi^n)\), then \(\varphi^n = 0\).
}

\nt{Over a non-algebraically closed field, (3) needs modification: \(\varphi\) is nilpotent iff \(\varphi\) has no eigenvalues in any field extension, or equivalently, iff the characteristic polynomial is \(x^n\).}

\cor{Nilpotent matrices are similar to strictly upper-triangular matrices}{If \(\varphi\) is nilpotent, then there exists a basis in which the matrix of \(\varphi\) is strictly upper-triangular.}

\pf{Proof}{We prove this by induction on \(n = \dim V\). If \(n = 1\), then \(\varphi = 0\) and we're done.

For \(n > 1\): since \(\varphi\) is nilpotent and nonzero (otherwise trivial), \(\Ker(\varphi) \neq \{0\}\). Pick \(0 \neq v_1 \in \Ker(\varphi)\) and extend to a basis \(\{v_1, v_2, \ldots, v_n\}\). The matrix of \(\varphi\) has first column all zeros. 

Alternatively, consider \(W = \Img(\varphi)\). Since \(\varphi\) is nilpotent and nonzero, \(W \subsetneq V\) (otherwise \(\varphi\) would be surjective, hence bijective, contradicting nilpotency). The restriction \(\varphi|_W : W \to W\) is well-defined (\(W\) is \(\varphi\)-invariant) and nilpotent. By induction, choose a basis of \(W\) making \(\varphi|_W\) strictly upper-triangular, extend to a basis of \(V\), and the result follows.
}


\subsection{The Nilpotent Basis Construction}

The key to understanding nilpotent operators is building a basis adapted to the ``chains'' of vectors connected by \(\varphi\).

\dfn{Chains and Jordan chains}{Let \(\varphi \in \End(V)\) be nilpotent. A \textbf{chain} (or \textbf{Jordan chain}) of length \(m\) is a sequence of vectors \(v_1, v_2, \ldots, v_m\) such that:
\[\varphi(v_m) = v_{m-1}, \quad \varphi(v_{m-1}) = v_{m-2}, \quad \ldots, \quad \varphi(v_2) = v_1, \quad \varphi(v_1) = 0.\]
The vector \(v_m\) is called the \textbf{head} of the chain, and \(v_1\) is the \textbf{tail}.}

\nt{In a chain, \(v_1 \in \Ker(\varphi)\), \(v_2 \in \Ker(\varphi^2) \setminus \Ker(\varphi)\), and generally \(v_j \in \Ker(\varphi^j) \setminus \Ker(\varphi^{j-1})\). The chain encodes how \(\varphi\) ``moves'' vectors toward the kernel.}

\mlemma{Key lemma for chain construction}{Let \(\varphi\) be nilpotent. Suppose \(U \subseteq \Ker(\varphi^{k+1})\) is a subspace satisfying \(U \cap \Ker(\varphi^k) = \{0\}\). Then:
\begin{enumerate}
    \item \(\varphi|_U\) is injective.
    \item \(\varphi(U) \subseteq \Ker(\varphi^k)\).
    \item \(\varphi(U) \cap \Ker(\varphi^{k-1}) = \{0\}\) (if \(k \geq 1\)).
\end{enumerate}
}

\pf{Proof}{(1) If \(v \in U\) and \(\varphi(v) = 0\), then \(v \in U \cap \Ker(\varphi) \subseteq U \cap \Ker(\varphi^k) = \{0\}\).

(2) For \(v \in U \subseteq \Ker(\varphi^{k+1})\), we have \(\varphi^{k+1}(v) = 0\), so \(\varphi^k(\varphi(v)) = 0\), meaning \(\varphi(v) \in \Ker(\varphi^k)\).

(3) If \(\varphi(v) \in \Ker(\varphi^{k-1})\), then \(\varphi^k(v) = \varphi^{k-1}(\varphi(v)) = 0\), so \(v \in \Ker(\varphi^k)\). By assumption \(U \cap \Ker(\varphi^k) = \{0\}\), so \(v = 0\), hence \(\varphi(v) = 0\).
}

\thm{Structure theorem for nilpotent operators}{Let \(\varphi \in \End(V)\) be nilpotent with index \(m\) (so \(\varphi^m = 0\) but \(\varphi^{m-1} \neq 0\)). Then there exists a basis of \(V\) of the form
\[\{v_{1,1}, \ldots, v_{1,m_1}; \; v_{2,1}, \ldots, v_{2,m_2}; \; \ldots; \; v_{\ell,1}, \ldots, v_{\ell, m_\ell}\}\]
where each sequence \((v_{i,1}, \ldots, v_{i,m_i})\) is a chain:
\[\varphi(v_{i,j}) = v_{i,j-1} \text{ for } j > 1, \quad \varphi(v_{i,1}) = 0.\]
The chain lengths satisfy \(m = m_1 \geq m_2 \geq \cdots \geq m_\ell \geq 1\) and \(\sum_{i=1}^\ell m_i = \dim V\).}

\pf{Proof}{We construct the basis by working down from \(\Ker(\varphi^m) = V\).

\textbf{Step 1:} Let \(U_m\) be a complement of \(\Ker(\varphi^{m-1})\) in \(V = \Ker(\varphi^m)\). Pick a basis \(\{v_{1,m}, \ldots, v_{k_m, m}\}\) of \(U_m\). These will be the heads of chains of length \(m\).

By the lemma, \(\varphi(U_m) \subseteq \Ker(\varphi^{m-1})\) and \(\varphi(U_m) \cap \Ker(\varphi^{m-2}) = \{0\}\). Define \(v_{i, m-1} = \varphi(v_{i,m})\) for each \(i\). These vectors are linearly independent (since \(\varphi|_{U_m}\) is injective).

\textbf{Step 2:} Now work in \(\Ker(\varphi^{m-1})\). We have \(\varphi(U_m) \subseteq \Ker(\varphi^{m-1})\) with \(\varphi(U_m) \cap \Ker(\varphi^{m-2}) = \{0\}\). Extend \(\varphi(U_m)\) to a larger subspace \(U_{m-1}\) that complements \(\Ker(\varphi^{m-2})\) in \(\Ker(\varphi^{m-1})\). Pick additional basis vectors \(\{v_{k_m+1, m-1}, \ldots, v_{k_{m-1}, m-1}\}\) to span the complement of \(\varphi(U_m)\) in \(U_{m-1}\). These start new chains of length \(m-1\).

\textbf{Continue inductively:} At each stage, we have chains of various lengths. Applying \(\varphi\) extends existing chains, and we add new chain heads as needed to complement \(\Ker(\varphi^{j-1})\) in \(\Ker(\varphi^j)\).

\textbf{Termination:} At the final step, we're in \(\Ker(\varphi)\), and all chains terminate. The collection of all chain vectors forms a basis of \(V\).
}

\nt{The chain lengths \((m_1, m_2, \ldots, m_\ell)\) form a partition of \(n = \dim V\). This partition is an invariant of the nilpotent operator---two nilpotent operators are similar iff they have the same partition.}


\subsection{Jordan Blocks for Nilpotent Operators}

The chain basis gives the matrix of a nilpotent operator a very specific form.

\dfn{Nilpotent Jordan block}{A \textbf{nilpotent Jordan block} of size \(m\) is the \(m \times m\) matrix
\[J_m(0) = \begin{pmatrix} 0 & 1 & 0 & \cdots & 0 \\ 0 & 0 & 1 & \cdots & 0 \\ \vdots & & \ddots & \ddots & \vdots \\ 0 & 0 & \cdots & 0 & 1 \\ 0 & 0 & \cdots & 0 & 0 \end{pmatrix}.\]
This has 1's on the superdiagonal and 0's elsewhere.}

\cor{Matrix form of nilpotent operators}{With respect to a chain basis ordered as \((v_{1,1}, \ldots, v_{1,m_1}, v_{2,1}, \ldots, v_{2,m_2}, \ldots)\), the matrix of a nilpotent operator \(\varphi\) is block diagonal:
\[\mcM(\varphi) = \begin{pmatrix} J_{m_1}(0) & & 0 \\ & J_{m_2}(0) & \\ 0 & & \ddots \end{pmatrix}.\]
}

\pf{Proof}{Within each chain \((v_{i,1}, \ldots, v_{i,m_i})\), we have \(\varphi(v_{i,j}) = v_{i,j-1}\) for \(j > 1\) and \(\varphi(v_{i,1}) = 0\). In matrix form with this basis ordering, this gives exactly \(J_{m_i}(0)\) for the \(i\)-th block. The blocks are independent since chains span complementary subspaces.
}

\ex{Nilpotent operator with partition \((3, 2, 1)\)}{Consider \(\varphi : k^6 \to k^6\) nilpotent with chain lengths \(3, 2, 1\). A chain basis might be:
\[\{v_{1,1}, v_{1,2}, v_{1,3}, v_{2,1}, v_{2,2}, v_{3,1}\}\]
where \(\varphi(v_{1,3}) = v_{1,2}\), \(\varphi(v_{1,2}) = v_{1,1}\), \(\varphi(v_{1,1}) = 0\), etc.

The matrix is:
\[\mcM(\varphi) = \begin{pmatrix} 0 & 1 & 0 & & & \\ 0 & 0 & 1 & & 0 & \\ 0 & 0 & 0 & & & \\ & & & 0 & 1 & \\ & 0 & & 0 & 0 & \\ & & & & & 0 \end{pmatrix} = J_3(0) \oplus J_2(0) \oplus J_1(0).\]
}

\mprop{Uniqueness of the Jordan form for nilpotent operators}{The partition \((m_1, m_2, \ldots, m_\ell)\) associated to a nilpotent operator \(\varphi\) is uniquely determined by \(\varphi\). Equivalently, two nilpotent operators are similar iff they have the same Jordan block sizes (counting multiplicity).}

\pf{Proof}{The partition can be recovered from dimensions of iterated kernels. Let \(d_j = \dim \Ker(\varphi^j)\). Then:
\[d_1 - d_0 = d_1 = \text{number of chains} = \ell\]
\[d_2 - d_1 = \text{number of chains of length} \geq 2\]
\[d_j - d_{j-1} = \text{number of chains of length} \geq j.\]

From this sequence of differences, we can reconstruct the partition: the number of blocks of size exactly \(j\) is \((d_j - d_{j-1}) - (d_{j+1} - d_j)\). Since the \(d_j\) are determined by \(\varphi\), so is the partition.
}

\nt{This dimension-counting argument is the key to proving uniqueness of Jordan form in general. The rank of \((\varphi - \lambda I)^j\) determines how many Jordan blocks for \(\lambda\) have size \(\geq j\).}


\section{Jordan Normal Form}

We now combine the generalized eigenspace decomposition with the structure theory of nilpotent operators to obtain the Jordan normal form---the ``best possible'' matrix representation for an operator over an algebraically closed field.

\subsection{Jordan Blocks}

\dfn{Jordan block}{A \textbf{Jordan block} of size \(m\) with eigenvalue \(\lambda\) is the \(m \times m\) matrix
\[J_m(\lambda) = \begin{pmatrix} \lambda & 1 & 0 & \cdots & 0 \\ 0 & \lambda & 1 & \cdots & 0 \\ \vdots & & \ddots & \ddots & \vdots \\ 0 & 0 & \cdots & \lambda & 1 \\ 0 & 0 & \cdots & 0 & \lambda \end{pmatrix} = \lambda I_m + J_m(0).\]
This has \(\lambda\) on the diagonal, 1's on the superdiagonal, and 0's elsewhere.}

\ex{Small Jordan blocks}{
\[J_1(\lambda) = (\lambda), \quad J_2(\lambda) = \begin{pmatrix} \lambda & 1 \\ 0 & \lambda \end{pmatrix}, \quad J_3(\lambda) = \begin{pmatrix} \lambda & 1 & 0 \\ 0 & \lambda & 1 \\ 0 & 0 & \lambda \end{pmatrix}.\]
A \(1 \times 1\) Jordan block is just a scalar---this corresponds to a true eigenvector. Larger blocks correspond to generalized eigenvectors that are not true eigenvectors.
}

\mprop{Properties of Jordan blocks}{Let \(J = J_m(\lambda)\). Then:
\begin{enumerate}
    \item The only eigenvalue of \(J\) is \(\lambda\), with algebraic multiplicity \(m\).
    \item The eigenspace \(E(\lambda, J) = \Ker(J - \lambda I) = \operatorname{Span}\{e_1\}\) has dimension 1.
    \item The generalized eigenspace is all of \(k^m\).
    \item \(J - \lambda I = J_m(0)\) is nilpotent with index \(m\).
\end{enumerate}
}

\pf{Proof}{(1) The characteristic polynomial is \(\det(J - xI) = (\lambda - x)^m\).

(2) \(J - \lambda I = J_m(0)\) has kernel spanned by \(e_1\) (the first standard basis vector).

(3) Since \((J - \lambda I)^m = J_m(0)^m = 0\), every vector is in \(\Ker(J - \lambda I)^m\).

(4) \(J_m(0)^{m-1} \neq 0\) (it maps \(e_m \mapsto e_1\)) but \(J_m(0)^m = 0\).
}

\nt{A Jordan block \(J_m(\lambda)\) represents the action of an operator on a single chain of generalized eigenvectors. The chain is \((e_1, e_2, \ldots, e_m)\) where \(Je_j = \lambda e_j + e_{j-1}\) for \(j > 1\) and \(Je_1 = \lambda e_1\). Equivalently, \((J - \lambda I)e_j = e_{j-1}\) for \(j > 1\) and \((J - \lambda I)e_1 = 0\).}


\subsection{The Jordan Normal Form Theorem}

\thm{Jordan normal form}{Let \(V\) be a finite-dimensional vector space over an algebraically closed field \(k\), and let \(\varphi \in \End(V)\). Then there exists a basis of \(V\) with respect to which the matrix of \(\varphi\) is block diagonal:
\[\mcM(\varphi) = \begin{pmatrix} J_{m_1}(\lambda_1) & & 0 \\ & J_{m_2}(\lambda_2) & \\ 0 & & \ddots \end{pmatrix}\]
where each block is a Jordan block. This matrix is called the \textbf{Jordan normal form} (or \textbf{Jordan canonical form}) of \(\varphi\).}

\pf{Proof}{By the generalized eigenspace decomposition, \(V = V_{\lambda_1} \oplus \cdots \oplus V_{\lambda_r}\) where \(\lambda_1, \ldots, \lambda_r\) are the distinct eigenvalues. Each \(V_{\lambda_i}\) is \(\varphi\)-invariant, so we can work on each generalized eigenspace separately.

On \(V_{\lambda_i}\), the operator \(\psi_i = \varphi|_{V_{\lambda_i}} - \lambda_i \id\) is nilpotent (this was proved earlier). By the structure theorem for nilpotent operators, there exists a chain basis of \(V_{\lambda_i}\) such that \(\psi_i\) has matrix \(J_{m_1}(0) \oplus J_{m_2}(0) \oplus \cdots\).

Since \(\varphi|_{V_{\lambda_i}} = \psi_i + \lambda_i \id\), the matrix of \(\varphi|_{V_{\lambda_i}}\) in this basis is
\[J_{m_1}(0) + \lambda_i I \oplus J_{m_2}(0) + \lambda_i I \oplus \cdots = J_{m_1}(\lambda_i) \oplus J_{m_2}(\lambda_i) \oplus \cdots\]

Combining the bases for all generalized eigenspaces gives a basis of \(V\) in which \(\varphi\) has the desired block diagonal form.
}

\nt{The Jordan form groups blocks by eigenvalue: all \(J_{m}(\lambda_i)\) blocks for a given \(\lambda_i\) appear together, corresponding to the decomposition of \(V_{\lambda_i}\) into chains. The total size of blocks with eigenvalue \(\lambda_i\) equals \(\dim V_{\lambda_i}\) = algebraic multiplicity of \(\lambda_i\).}


\subsection{Existence and Uniqueness}

\thm{Uniqueness of Jordan normal form}{The Jordan normal form of \(\varphi\) is unique up to permutation of the blocks. That is, the multiset of Jordan blocks \(\{J_{m_1}(\lambda_1), J_{m_2}(\lambda_2), \ldots\}\) is uniquely determined by \(\varphi\).}

\pf{Proof}{For each eigenvalue \(\lambda\), the Jordan blocks with eigenvalue \(\lambda\) correspond to the partition of \(\dim V_\lambda\) given by the nilpotent operator \(\varphi|_{V_\lambda} - \lambda \id\). We showed this partition is uniquely determined by dimensions of iterated kernels:
\[d_j(\lambda) = \dim \Ker((\varphi - \lambda \id)^j).\]

The number of blocks \(J_m(\lambda)\) of size exactly \(m\) is
\[(d_m(\lambda) - d_{m-1}(\lambda)) - (d_{m+1}(\lambda) - d_m(\lambda)) = 2d_m(\lambda) - d_{m-1}(\lambda) - d_{m+1}(\lambda).\]

Since these dimensions are intrinsic to \(\varphi\), the block structure is uniquely determined.
}

\cor{Similarity criterion}{Two operators \(\varphi, \psi \in \End(V)\) are similar if and only if they have the same Jordan normal form (up to block ordering). Equivalently, \(\varphi \sim \psi\) iff for each \(\lambda \in k\) and each \(j \geq 1\):
\[\dim \Ker((\varphi - \lambda \id)^j) = \dim \Ker((\psi - \lambda \id)^j).\]}

\nt{This gives a complete set of similarity invariants. Unlike trace and determinant alone (which don't distinguish all similarity classes), the full sequence of kernel dimensions completely determines the similarity class.}


\subsection{Computing Jordan Form: The Algorithm}

Given a matrix \(A\), here is the procedure to find its Jordan form:

\textbf{Step 1: Find eigenvalues.} Compute \(\chi_A(x) = \det(xI - A)\) and factor it as \(\prod_i (x - \lambda_i)^{n_i}\). The \(\lambda_i\) are eigenvalues with algebraic multiplicities \(n_i\).

\textbf{Step 2: For each eigenvalue \(\lambda\), determine block sizes.} Compute \(d_j = \dim \Ker(A - \lambda I)^j\) for \(j = 1, 2, \ldots\) until stabilization. The number of blocks of size \(\geq j\) is \(d_j - d_{j-1}\). Alternatively:
\begin{itemize}
    \item Number of blocks = \(d_1\) = geometric multiplicity.
    \item Number of blocks of size \(\geq 2\) is \(d_2 - d_1\).
    \item Number of blocks of size exactly \(j\) is \((d_j - d_{j-1}) - (d_{j+1} - d_j)\).
\end{itemize}

\textbf{Step 3: Assemble the Jordan form.} List all Jordan blocks grouped by eigenvalue.

\ex{Computing Jordan form}{Let \(A = \begin{pmatrix} 2 & 1 & 0 \\ 0 & 2 & 1 \\ 0 & 0 & 2 \end{pmatrix}\).

\textbf{Step 1:} \(\chi_A(x) = (x-2)^3\), so \(\lambda = 2\) with multiplicity 3.

\textbf{Step 2:} 
\[A - 2I = \begin{pmatrix} 0 & 1 & 0 \\ 0 & 0 & 1 \\ 0 & 0 & 0 \end{pmatrix}, \quad (A-2I)^2 = \begin{pmatrix} 0 & 0 & 1 \\ 0 & 0 & 0 \\ 0 & 0 & 0 \end{pmatrix}, \quad (A-2I)^3 = 0.\]
So \(d_1 = 1\), \(d_2 = 2\), \(d_3 = 3\).

Number of blocks = \(d_1 = 1\). Size of the single block = 3.

\textbf{Step 3:} Jordan form is \(J_3(2)\). (The matrix is already in Jordan form!)
}

\ex{A more interesting example}{Let \(A = \begin{pmatrix} 3 & 1 & 0 & 0 \\ 0 & 3 & 0 & 0 \\ 0 & 0 & 3 & 1 \\ 0 & 0 & 0 & 3 \end{pmatrix}\).

\textbf{Step 1:} \(\chi_A(x) = (x-3)^4\), so \(\lambda = 3\) with multiplicity 4.

\textbf{Step 2:}
\[A - 3I = \begin{pmatrix} 0 & 1 & 0 & 0 \\ 0 & 0 & 0 & 0 \\ 0 & 0 & 0 & 1 \\ 0 & 0 & 0 & 0 \end{pmatrix}.\]
\(\Ker(A - 3I) = \operatorname{Span}\{e_1, e_3\}\), so \(d_1 = 2\).
\((A - 3I)^2 = 0\), so \(d_2 = 4\).

Number of blocks = 2. Blocks of size \(\geq 2\): \(d_2 - d_1 = 2\). So both blocks have size 2.

\textbf{Step 3:} Jordan form is \(J_2(3) \oplus J_2(3)\).
}

\ex{Multiple eigenvalues}{Let \(A = \begin{pmatrix} 1 & 1 & 0 & 0 \\ 0 & 1 & 0 & 0 \\ 0 & 0 & 2 & 0 \\ 0 & 0 & 0 & 2 \end{pmatrix}\).

\textbf{Step 1:} \(\chi_A(x) = (x-1)^2(x-2)^2\). Eigenvalues: \(\lambda_1 = 1\), \(\lambda_2 = 2\), each with multiplicity 2.

\textbf{Step 2 for \(\lambda = 1\):}
\[A - I = \begin{pmatrix} 0 & 1 & 0 & 0 \\ 0 & 0 & 0 & 0 \\ 0 & 0 & 1 & 0 \\ 0 & 0 & 0 & 1 \end{pmatrix}.\]
Restricted to generalized eigenspace (first two coordinates): \(d_1 = 1\), \(d_2 = 2\). One block of size 2: \(J_2(1)\).

\textbf{Step 2 for \(\lambda = 2\):}
\[A - 2I = \begin{pmatrix} -1 & 1 & 0 & 0 \\ 0 & -1 & 0 & 0 \\ 0 & 0 & 0 & 0 \\ 0 & 0 & 0 & 0 \end{pmatrix}.\]
Restricted to generalized eigenspace (last two coordinates): \(d_1 = 2\). Two blocks of size 1: \(J_1(2) \oplus J_1(2)\).

\textbf{Step 3:} Jordan form is \(J_2(1) \oplus J_1(2) \oplus J_1(2) = \begin{pmatrix} 1 & 1 & 0 & 0 \\ 0 & 1 & 0 & 0 \\ 0 & 0 & 2 & 0 \\ 0 & 0 & 0 & 2 \end{pmatrix}\).
}

\nt{The last example shows that \(\lambda = 2\) contributes two \(1 \times 1\) blocks because the eigenspace has dimension 2 (geometric = algebraic multiplicity). This is the diagonalizable case for that eigenvalue. In contrast, \(\lambda = 1\) has a \(2 \times 2\) block because geometric multiplicity (1) is less than algebraic multiplicity (2).}


\section{Characteristic and Minimal Polynomials}

We now introduce two polynomials canonically associated to a linear operator. The characteristic polynomial encodes eigenvalues and their algebraic multiplicities; the minimal polynomial captures the ``smallest'' polynomial relation satisfied by the operator.

\subsection{The Characteristic Polynomial}

\dfn{Characteristic polynomial}{Let \(\varphi \in \End(V)\) with \(\dim V = n\), and let \(A = \mcM(\varphi)\) be the matrix of \(\varphi\) with respect to some basis. The \textbf{characteristic polynomial} of \(\varphi\) is
\[\chi_\varphi(x) = \det(xI - \varphi) = \det(xI - A).\]
This is a monic polynomial of degree \(n\).}

\mprop{Basic properties of the characteristic polynomial}{
\begin{enumerate}
    \item \(\chi_\varphi(x)\) is independent of the choice of basis (hence well-defined for the operator \(\varphi\)).
    \item \(\chi_\varphi(x) = x^n - (\Tr \varphi)x^{n-1} + \cdots + (-1)^n \det(\varphi)\).
    \item The roots of \(\chi_\varphi(x)\) are exactly the eigenvalues of \(\varphi\).
    \item If \(k\) is algebraically closed, \(\chi_\varphi(x) = \prod_{i=1}^m (x - \lambda_i)^{n_i}\) where \(\lambda_i\) are the distinct eigenvalues and \(n_i = \dim V_{\lambda_i}\).
\end{enumerate}
}

\pf{Proof}{(1) If \(B = P^{-1}AP\) is the matrix in another basis, then
\[\det(xI - B) = \det(xI - P^{-1}AP) = \det(P^{-1}(xI - A)P) = \det(xI - A).\]

(2) Expand \(\det(xI - A)\) using the Leibniz formula. The leading term comes from the diagonal product \(\prod_i (x - a_{ii}) = x^n - (\sum_i a_{ii})x^{n-1} + \cdots\). The constant term is \(\det(-A) = (-1)^n \det(A)\).

(3) \(\lambda\) is a root iff \(\det(\lambda I - A) = 0\) iff \(\lambda I - A\) is singular iff \(\Ker(\varphi - \lambda \id) \neq \{0\}\) iff \(\lambda\) is an eigenvalue.

(4) This follows from the generalized eigenspace decomposition: the algebraic multiplicity of \(\lambda_i\) equals \(\dim V_{\lambda_i}\).
}

\ex{Characteristic polynomial of a \(2 \times 2\) matrix}{For \(A = \begin{pmatrix} a & b \\ c & d \end{pmatrix}\):
\[\chi_A(x) = \det\begin{pmatrix} x - a & -b \\ -c & x - d \end{pmatrix} = (x-a)(x-d) - bc = x^2 - (a+d)x + (ad - bc) = x^2 - (\Tr A)x + \det A.\]
}


\subsection{The Cayley-Hamilton Theorem}

The Cayley-Hamilton theorem states that every operator satisfies its own characteristic polynomial.

\thm{Cayley-Hamilton}{Let \(\varphi \in \End(V)\) with characteristic polynomial \(\chi_\varphi(x)\). Then
\[\chi_\varphi(\varphi) = 0.\]
That is, substituting \(\varphi\) for \(x\) in \(\chi_\varphi(x)\) gives the zero operator.}

\pf{Proof}{We give two proofs.

\textbf{Proof 1 (via Jordan form, over algebraically closed \(k\)):} It suffices to prove this for Jordan blocks, since the Jordan form is block diagonal and \(\chi_\varphi = \prod \chi_{J_i}\).

For a single Jordan block \(J = J_m(\lambda)\), we have \(\chi_J(x) = (x - \lambda)^m\). Then
\[\chi_J(J) = (J - \lambda I)^m = J_m(0)^m = 0\]
since \(J_m(0)\) is nilpotent with index \(m\).

\textbf{Proof 2 (general, using adjugate matrix):} Let \(A\) be the matrix of \(\varphi\). The adjugate (classical adjoint) matrix \(\operatorname{adj}(xI - A)\) is a matrix whose entries are polynomials in \(x\) of degree at most \(n-1\). Write
\[\operatorname{adj}(xI - A) = B_{n-1}x^{n-1} + B_{n-2}x^{n-2} + \cdots + B_1 x + B_0\]
where \(B_i\) are constant matrices.

By the adjugate identity: \((xI - A) \cdot \operatorname{adj}(xI - A) = \det(xI - A) \cdot I = \chi_A(x) \cdot I\).

Expanding both sides and comparing coefficients of powers of \(x\) gives a system of matrix equations. Multiplying the equation for \(x^k\) by \(A^k\) and summing yields \(\chi_A(A) = 0\). (The details involve telescoping---each \(B_i\) appears twice with opposite signs.)
}

\nt{The Cayley-Hamilton theorem does \emph{not} follow from ``plugging in \(A\) for \(x\) in \(\det(xI - A) = 0\).'' The determinant is a scalar-valued function, not a matrix-valued function, so this substitution is meaningless. The theorem requires proof.}

\cor{Bound on nilpotence index}{If \(\varphi\) is nilpotent on an \(n\)-dimensional space, then \(\varphi^n = 0\).}

\pf{Proof}{The only eigenvalue is 0, so \(\chi_\varphi(x) = x^n\). By Cayley-Hamilton, \(\varphi^n = 0\).
}

\cor{Expression for the inverse}{If \(\varphi\) is invertible, then \(\varphi^{-1}\) can be expressed as a polynomial in \(\varphi\).}

\pf{Proof}{Write \(\chi_\varphi(x) = x^n + c_{n-1}x^{n-1} + \cdots + c_1 x + c_0\). Since \(\varphi\) is invertible, \(0\) is not an eigenvalue, so \(c_0 = (-1)^n \det(\varphi) \neq 0\).

By Cayley-Hamilton: \(\varphi^n + c_{n-1}\varphi^{n-1} + \cdots + c_1 \varphi + c_0 I = 0\).

Rearranging: \(\varphi(\varphi^{n-1} + c_{n-1}\varphi^{n-2} + \cdots + c_1 I) = -c_0 I\).

Thus \(\varphi^{-1} = -\frac{1}{c_0}(\varphi^{n-1} + c_{n-1}\varphi^{n-2} + \cdots + c_1 I)\).
}


\subsection{The Minimal Polynomial}

\dfn{Minimal polynomial}{The \textbf{minimal polynomial} \(\mu_\varphi(x)\) of \(\varphi \in \End(V)\) is the unique monic polynomial of smallest degree such that \(\mu_\varphi(\varphi) = 0\).}

\mprop{Existence and uniqueness of the minimal polynomial}{
\begin{enumerate}
    \item There exists a nonzero polynomial \(p(x)\) with \(p(\varphi) = 0\).
    \item Among all such polynomials, there is a unique monic one of minimal degree.
    \item \(\mu_\varphi(x)\) divides every polynomial \(p(x)\) satisfying \(p(\varphi) = 0\).
\end{enumerate}
}

\pf{Proof}{(1) Cayley-Hamilton gives \(\chi_\varphi(\varphi) = 0\).

(2) Let \(d\) be the minimal degree of a nonzero polynomial annihilating \(\varphi\). If \(p(x)\) and \(q(x)\) are both monic of degree \(d\) with \(p(\varphi) = q(\varphi) = 0\), then \((p - q)(\varphi) = 0\) and \(\deg(p - q) < d\). By minimality, \(p - q = 0\).

(3) Given \(p(x)\) with \(p(\varphi) = 0\), perform polynomial division: \(p(x) = q(x)\mu_\varphi(x) + r(x)\) where \(\deg r < \deg \mu_\varphi\). Then \(r(\varphi) = p(\varphi) - q(\varphi)\mu_\varphi(\varphi) = 0 - 0 = 0\). By minimality of \(\mu_\varphi\), we must have \(r = 0\), so \(\mu_\varphi \mid p\).
}

\thm{Relationship between characteristic and minimal polynomials}{Let \(\varphi \in \End(V)\).
\begin{enumerate}
    \item \(\mu_\varphi(x)\) divides \(\chi_\varphi(x)\).
    \item \(\mu_\varphi(x)\) and \(\chi_\varphi(x)\) have the same roots (i.e., the same irreducible factors over \(k\)).
    \item Over an algebraically closed field, if \(\chi_\varphi(x) = \prod_i (x - \lambda_i)^{n_i}\), then \(\mu_\varphi(x) = \prod_i (x - \lambda_i)^{m_i}\) where \(1 \leq m_i \leq n_i\).
\end{enumerate}
}

\pf{Proof}{(1) By Cayley-Hamilton, \(\chi_\varphi(\varphi) = 0\). By the divisibility property of \(\mu_\varphi\), we have \(\mu_\varphi \mid \chi_\varphi\).

(2) If \(\lambda\) is a root of \(\chi_\varphi\), then \(\lambda\) is an eigenvalue: there exists \(v \neq 0\) with \(\varphi(v) = \lambda v\). For any polynomial \(p\), \(p(\varphi)(v) = p(\lambda)v\). If \(p(\varphi) = 0\), then \(p(\lambda) = 0\). In particular, \(\mu_\varphi(\lambda) = 0\).

Conversely, if \(\mu_\varphi(\lambda) = 0\), then \((x - \lambda) \mid \mu_\varphi(x)\), so \((x - \lambda) \mid \chi_\varphi(x)\) by (1), meaning \(\lambda\) is a root of \(\chi_\varphi\).

(3) From (1) and (2), the factorizations have the same roots with \(m_i \leq n_i\). The exponent \(m_i\) is the size of the largest Jordan block with eigenvalue \(\lambda_i\) (proof below).
}

\mprop{Minimal polynomial and Jordan block sizes}{The exponent of \((x - \lambda)\) in \(\mu_\varphi(x)\) equals the size of the largest Jordan block with eigenvalue \(\lambda\).}

\pf{Proof}{On the generalized eigenspace \(V_\lambda\), the operator \(\varphi - \lambda \id\) is nilpotent. Its nilpotence index is the size of the largest Jordan block (the longest chain). Call this \(m_\lambda\).

For \((\varphi - \lambda \id)^k\) to annihilate \(V_\lambda\), we need \(k \geq m_\lambda\). The smallest such \(k\) is \(m_\lambda\).

Since \(V = \bigoplus_\lambda V_\lambda\) and these are \(\varphi\)-invariant, \(p(\varphi) = 0\) iff \(p(\varphi)|_{V_\lambda} = 0\) for all \(\lambda\). The latter requires \((x - \lambda)^{m_\lambda} \mid p(x)\). Thus \(\mu_\varphi(x) = \prod_\lambda (x - \lambda)^{m_\lambda}\).
}


\subsection{Diagonalizability Criterion}

\thm{Diagonalizability via minimal polynomial}{An operator \(\varphi \in \End(V)\) is diagonalizable if and only if \(\mu_\varphi(x)\) splits into distinct linear factors:
\[\mu_\varphi(x) = (x - \lambda_1)(x - \lambda_2) \cdots (x - \lambda_m)\]
where \(\lambda_1, \ldots, \lambda_m\) are distinct.}

\pf{Proof}{(\(\Rightarrow\)) If \(\varphi\) is diagonalizable, all Jordan blocks are \(1 \times 1\). The largest block for each eigenvalue has size 1, so each \((x - \lambda_i)\) appears with exponent 1 in \(\mu_\varphi\).

(\(\Leftarrow\)) If \(\mu_\varphi(x) = \prod_i (x - \lambda_i)\) with distinct \(\lambda_i\), then each Jordan block has size 1 (since the exponent of \((x - \lambda_i)\) in \(\mu_\varphi\) equals the largest block size for \(\lambda_i\)). So \(\varphi\) has a basis of eigenvectors.
}

\cor{Diagonalizability of operators with distinct eigenvalues}{If \(\varphi\) has \(n = \dim V\) distinct eigenvalues, then \(\varphi\) is diagonalizable.}

\pf{Proof}{The characteristic polynomial has \(n\) distinct roots, so \(\chi_\varphi(x) = \prod_{i=1}^n (x - \lambda_i)\). Since \(\mu_\varphi \mid \chi_\varphi\) and they have the same roots, \(\mu_\varphi\) also has distinct linear factors.
}

\ex{Minimal polynomial examples}{
\begin{enumerate}
    \item For \(A = \begin{pmatrix} 2 & 0 \\ 0 & 2 \end{pmatrix} = 2I\): \(\chi_A(x) = (x-2)^2\), but \(\mu_A(x) = x - 2\) since \((A - 2I) = 0\).
    
    \item For \(A = \begin{pmatrix} 2 & 1 \\ 0 & 2 \end{pmatrix}\): \(\chi_A(x) = (x-2)^2\). Since \(A - 2I = \begin{pmatrix} 0 & 1 \\ 0 & 0 \end{pmatrix} \neq 0\), we need \(\mu_A(x) = (x-2)^2 = \chi_A(x)\).
    
    \item For \(A = \diag(1, 1, 2, 2, 2)\): \(\chi_A(x) = (x-1)^2(x-2)^3\), but \(\mu_A(x) = (x-1)(x-2)\) since \((A-I)(A-2I) = 0\).
\end{enumerate}
}


\subsection{Exercises}

\qs{Exercises on advanced spectral theory}{
\begin{enumerate}
    \item Let \(\varphi : \CC^4 \to \CC^4\) have characteristic polynomial \(\chi_\varphi(x) = (x-1)^2(x-2)^2\). List all possible Jordan normal forms for \(\varphi\).
    
    \item Prove that if \(A\) and \(B\) are similar matrices, then \(\mu_A = \mu_B\).
    
    \item Find the Jordan normal form and minimal polynomial of \(A = \begin{pmatrix} 0 & 1 & 0 \\ 0 & 0 & 1 \\ 0 & 0 & 0 \end{pmatrix}\).
    
    \item Let \(\varphi \in \End(V)\) satisfy \(\varphi^2 = \varphi\). Prove that \(\varphi\) is diagonalizable and find its possible eigenvalues.
    
    \item Suppose \(\varphi \in \End(V)\) has minimal polynomial \(\mu_\varphi(x) = (x-1)(x-2)(x-3)\). What can you conclude about \(\varphi\)?
\end{enumerate}
}

\begin{solution}
\textbf{(1)} The eigenvalues are \(\lambda_1 = 1\) and \(\lambda_2 = 2\), each with algebraic multiplicity 2. For each eigenvalue, the Jordan blocks must have sizes summing to 2. Possibilities for each:
\begin{itemize}
    \item One block of size 2: \(J_2(\lambda)\)
    \item Two blocks of size 1: \(J_1(\lambda) \oplus J_1(\lambda)\)
\end{itemize}

Combining independently for each eigenvalue, the possible Jordan forms are:
\begin{enumerate}[(i)]
    \item \(J_2(1) \oplus J_2(2)\)
    \item \(J_2(1) \oplus J_1(2) \oplus J_1(2)\)
    \item \(J_1(1) \oplus J_1(1) \oplus J_2(2)\)
    \item \(J_1(1) \oplus J_1(1) \oplus J_1(2) \oplus J_1(2)\) (diagonal matrix)
\end{enumerate}

\textbf{(2)} If \(B = P^{-1}AP\), then for any polynomial \(p\):
\[p(B) = p(P^{-1}AP) = P^{-1}p(A)P.\]
So \(p(B) = 0\) iff \(p(A) = 0\). The set of annihilating polynomials is the same, hence the minimal polynomial (the monic generator of this ideal) is the same.

\textbf{(3)} The matrix is already in Jordan form: \(A = J_3(0)\). 

Characteristic polynomial: \(\chi_A(x) = x^3\).

Minimal polynomial: We need the smallest \(k\) with \(A^k = 0\). Compute \(A^2 = \begin{pmatrix} 0 & 0 & 1 \\ 0 & 0 & 0 \\ 0 & 0 & 0 \end{pmatrix} \neq 0\) and \(A^3 = 0\). So \(\mu_A(x) = x^3 = \chi_A(x)\).

\textbf{(4)} The condition \(\varphi^2 = \varphi\) means \(\varphi^2 - \varphi = 0\), so \(\varphi(\varphi - \id) = 0\). Thus \(p(x) = x(x-1) = x^2 - x\) annihilates \(\varphi\).

The minimal polynomial divides \(x(x-1)\). Since \(x(x-1)\) has distinct roots, so does \(\mu_\varphi\). By the diagonalizability criterion, \(\varphi\) is diagonalizable.

The possible eigenvalues are roots of \(\mu_\varphi\), which divides \(x(x-1)\). So eigenvalues are among \(\{0, 1\}\).

(Such operators are called \textbf{projections} or \textbf{idempotents}. They decompose \(V = \Ker(\varphi) \oplus \Img(\varphi)\) and project onto \(\Img(\varphi)\) along \(\Ker(\varphi)\).)

\textbf{(5)} Since \(\mu_\varphi(x) = (x-1)(x-2)(x-3)\) has three distinct linear factors:
\begin{itemize}
    \item \(\varphi\) is diagonalizable (by the criterion).
    \item The eigenvalues are exactly \(1, 2, 3\) (roots of \(\mu_\varphi\)).
    \item \(\dim V \geq 3\) (at least one eigenvector per eigenvalue).
    \item \(V = E(1, \varphi) \oplus E(2, \varphi) \oplus E(3, \varphi)\).
\end{itemize}
The eigenspace dimensions can be any positive integers summing to \(\dim V\), but we need at least one of each.
\end{solution}


\end{document}
